% concatenated from 20251230_distance.tex
 \documentclass[10pt]{amsart}

\usepackage[latin1]{inputenc}
\usepackage[english]{babel}
\usepackage{amsmath}
\usepackage{amsfonts}
\usepackage{amssymb}
\usepackage{amsthm}
\usepackage{epsfig}
\usepackage{graphicx,colortbl}
\usepackage{url,hyperref}
\usepackage{mathtools, color, changes}
\usepackage{dsfont}

\newtheorem{theorem}{Theorem}[section]
\newtheorem{lemma}{Lemma}[section]
\newtheorem{cor}{Corollary}[section]
\newtheorem{proposition}{Proposition}[section]
\newtheorem{definition}{Definition}[section]
\newtheorem{conjecture}{Conjecture}[section]
\newtheorem{question}{Question}[section]
\newtheorem{remark}{Remark}[section]
\newtheorem{rmk}{Remark}[section]
\newtheorem{example}{Example}[section]

%\newtheorem{thm}{Theorem}[section]
%\newtheorem{defin}[thm]{Definition}
%\newtheorem{thantiguo}{Theorem}
%\newtheorem{cor}[thm]{Corollary}
%\newtheorem{lem}[thm]{Lemma}
%\newtheorem{prop}[thm]{Proposition}
%{ \theoremstyle{remark}\newtheorem*{remark}{Remark} }
%\newtheorem{rmk}[thm]{Remark}

\newcounter{theor}
\renewcommand{\thetheor}{\Alph{theor}}
\newtheorem{thm}[theor]{Theorem}

\DeclareMathOperator{\inter}{int}
\def\pos{\mathop\mathrm{pos}\nolimits}
\def\aff{\mathop\mathrm{aff}\nolimits}
\def\lin{\mathop\mathrm{lin}\nolimits}
\def\conv{\mathop\mathrm{conv}\nolimits}
\def\bd{\mathop\mathrm{bd}\nolimits}
\def\cl{\mathop\mathrm{cl}\nolimits}
\def\supp{\mathop\mathrm{supp}\nolimits}
\def\hyp{\mathop\mathrm{hyp}\nolimits}
\def\B{\mathbb{B}}
\def\s{\mathbb{S}}
\def\K{\mathcal{K}}
\def\w{\mathcal{W}}
\def\S{\mathrm{S}}
\def\L{\mathcal{L}}
\def\r{\mathcal{R}}
\def\rd{\widetilde{\mathcal{R}}}
\def\R{\mathbb{R}}
\def\N{\mathbb{N}}
\def\Z{\mathbb{Z}}
\def\Q{\mathbb{Q}}
\def\V{\mathrm{V}}
\def\vol{\mathrm{Vol}}
\def\F{\mathcal{F}}
\def\D{\mathrm{D}}
\def\C{\mathbb{C}}
\def\c{\mathcal{C}}
\def\M{\mathrm{M}}
\def\b{\mathrm{b}}
\def\W{\mathrm{W}}
\def\e{\mathrm{e}}
\newcommand{\dlat}{\mathrm{d}}
\newcommand{\enorm}[1]{\left|{#1}\right|}

\def\E{\mathbb{E}}

\def\Var{\text{Var}}
\def\var{\text{Var}}
\def\cov{\text{Cov}}
\def\dist{\text{dist}}

\def\e{\mathrm{e}}
\def\esc#1{\left\langle #1\right\rangle}
\def\C{\mathcal{C}}
\def\I{\mathcal{I}}
\def\Cube#1{\mathrm{C}_{#1}}
\def\d{\text{d}}
\def\ceil#1{\lceil#1\rceil}
\def\floor#1{\lfloor#1\rfloor}
\def\2Z{2^{-m}\Z^n}
%Funcionales
\def\Gsub#1{\mathrm{G}_{#1}}
\def\G{\mathrm{G}_n}
%extensi�n de funci�n
\def\symbol{\diamond}
%Medias y sumas
\def\Media#1#2#3#4{\mathcal{M}_{#1}^{#4}\left(#2,#3\right)}
\def\Suma#1#2#3#4{\mathcal{S}_{#1}^{#4}\left(#2,#3\right)}
%Otro conjuntos
\def\Rm{\mathrm{R}^m}
\def\Om{\mathrm{O}^m}
\DeclareMathOperator*{\esssup}{ess\,sup}

\numberwithin{equation}{section}
\renewcommand{\labelenumi}{\rm \roman{enumi})}

\begin{document}

\title[CLT for distance]{Limit theorems for the distance \\of random points in $l_p^n$-balls}
% or if you want, simply \title{Title of the article}
\author{David Alonso-Guti\'errez, Javier Mart\'in Go\~ni, Joscha Prochno}
\address{University of Passau. Faculty of Computer Sciences and Mathematics. Innstrasse 41, 94032 Passau, Germany.}
\email{j.martin@unizar.es, javier.martingoni@uni-passau.de}
\thanks{The first named author is partially supported by MICINN project PID2022-137294NB-I00 and DGA project E48\_23R. The second and thirsd authors are supported by the DFG project 516672205}
\maketitle

\begin{abstract}
In this paper, we prove that the Euclidean distance between two independent random vectors uniformly distributed on $l_p^n$-balls ($1 \leq p \leq \infty$) or on its boundary satisfies a central limit theorem as $n\to\infty$. Also, we give a compact proof of the case of the sphere, which was proved by Hammersley in \cite{H}. Furthermore, we complement our central limit theorem by providing large deviation principles for the cases $p \geq 2$.
\end{abstract}

%\begin{abstract}
%In this paper, we sudy the euclidean distance between two random vectors uniformly distributed on $l_p^n$-balls ($1 \leq p \leq \infty$) and on its boundary. We provide a Central Limit Theorem. Also, we give a compact proof of the case of the sphere, which was proved by Hammersley in (...). Furthermore, we obtain some Large Deviation Principles for the cases where $p \geq 2$ and for the hypercube.
%\end{abstract}




 \section{Introduction, notation and results.}
 \label{section1}

Recent developments in Asymptotic Geometric Analysis have significantly deepened our understanding of high-dimensional spaces. This interdisciplinary field, which merges techniques from analysis, geometry, and probability, has unveiled striking regularities and patterns in geometric structures as the dimension grows. In particular, the behavior of random objects in high dimensions has been a central focus, revealing connections to classical probabilistic phenomena.

One prominent area of study concerns central limit theorems (CLT) in high-dimensional geometry. Early works by Poincar\'e and Borel (see \cite{DF} or \cite{St} for modern expositions) demonstrated that as the dimension of Euclidean spaces increases, the distribution of random points in high-dimensional spheres or balls tends to converge to a Gaussian distribution. This convergence has since been extended to more general convex bodies, notably through Klartag's CLT \cite{K,K2}, which establishes that projections of random vectors on high-dimensional convex sets exhibit Gaussian-like behavior.

Further advances in the field have focused on understanding the probabilistic structure of high-dimensional bodies, such as cubes and  $l_p^n$-balls. The work of Paouris, Pivovarov, and Zinn \cite{PPZ} has contributed important results in this direction, providing a central limit theorem for the volume of $k$-dimensional random projections ($k$ fixed) of the $n$-dimensional cube as the space dimension $n$ tends to infinity. Prochno, Th\"ale and Tuchel \cite{PTT} generalized and complemented this considerably by establishing a central limit theorem and moderate and large deviation principles for the volume of random projections and sections of $l_p^n$-balls. Moreover, there is a central limit theorem for the volume of convex hulls of Gaussian random vectors obtained by B\'ar\'any and Vu \cite{VB} or Reitzner's central limit theorems for the volume and the number of $i$-dimensional faces of random polytopes in smooth convex bodies \cite{R}.

Beyond these normal fluctuations, there has been growing interest in understanding large deviations in high-dimensional spaces. These deviations capture rare events where the behavior of random vectors significantly differs from typical Gaussian outcomes, highlighting non-universal features of different convex bodies. In particular, the study of large deviation principles (LDPs) has provided critical insights into these rare events, offering a more complete picture of the probabilistic behavior of random vectors in convex bodies; we refer the interested reader to \cite{P} and the references cited therein. 


The aim of this paper is to investigate CLTs and LDPs related to the Euclidean distance between random vectors in $l_p^n$-balls, to uncover deeper connections between geometric properties and probabilistic phenomena in high-dimensional settings. This problem was first researched by Hammersley in \cite{H}, where he proved an initial CLT for the case of the hypersphere. Further, it was extended to spherical distributions by Lord in \cite{L}. In this paper, we extend those results to $l_p^n$-balls and their respective boundaries. Furthermore, we obtain LDPs for the regime $2\leq p \leq \infty$. 



%in the context of high-dimensional convex bodies, which was first investigated by (), further by (). In particular, we will be focusing on the distance between random vectors in $l_p$-balls. By extending known results in these areas, we hope to uncover deeper connections between geometric properties and probabilistic phenomena in high-dimensional settings. 
 
We start with the central limit theorems and then present the corresponding large deviations counterparts. For $p \in [1, \infty ]$ and $x = (x_1, \ldots , x_n) \in \R^n$, $n\in \N$, we recall the definition of the $\Vert \cdot \Vert _p$-norm as 
\[   
\Vert x \Vert _p := 
     \begin{cases}
       \left( \sum _{i=1} ^n |x_i|^p \right)^{1/p} &\quad : p < \infty\\
       \max _{1 \leq i \leq n} |x_i| &\quad : p = \infty . \\
     \end{cases}
\]
We denote the unit ball in $\R^n$ with respect to the $\Vert \cdot \Vert _p$-norm by $B_p^n$, that is, 
\begin{align*}
B_p^n := \lbrace x \in \R^n \; \Vert x \Vert _p \leq 1 \rbrace,
\end{align*}
and we write $\partial B_p^n$ for its boundary. We write $X \sim \text{Unif} ( B_p^n)$ or $X \sim \text{Unif} ( \partial B_p^n)$ to denote that $X$ is a random vector uniformly distributed on $B_p^n$ or on $ \partial B_p^n$, with respect to the cone measure (see Section \ref{Section2.2}), respectively.



\subsection{Central limit theorems}

A cornerstone in probability theory is the central limit theorem. A sequence $(X_n)_{n \in \N}$
of random variables is said to satisfy the central limit theorem if 
\begin{align*}
\frac{X_n - \mathbb{E} X_n}{\sqrt{\Var X_n}} \overset{d}{\rightarrow} N(0,1),
\end{align*} 
where $\overset{d}{\rightarrow}$ denotes convergence in distribution as $n\to\infty$. That is, for every $t \in \R$,
\begin{align*}
\lim _{n \to \infty} \mathbb{P} \left( \frac{X_n - \mathbb{E} X_n}{\sqrt{\Var X_n}} \leq t \right) - \Phi (t) = 0,
\end{align*}
where $\Phi (t)$ is the distribution function of the standard normal distribution.
More generally, we say that a sequence of random variables $(Y_n)_{n \in \N}$ satisfies a CLT if $Y_n \overset{d}{\to} N(0, \sigma ^2)$, for some $\sigma > 0$.


Let us introduce the followings quantities 
\begin{align*}
M_p(\alpha) := \frac{\Gamma \left( (  \alpha + 1) / p \right) }{ \Gamma (1 / p )} ,
\end{align*}
where $\alpha >0$ and $p \in [1, \infty ]$.

In this paper, we are going to prove the following theorems. This first theorem establishes a central limit theorem for the distance between $X^{(n)}$ and $Y^{(n)}$, which are random vectors uniformly distributed in $l_p^n$-balls or on their boundary.

\begin{theorem}
\label{Thm:CLTforlpballs}
Let $1\leq p <\infty$, and let $X^{(n)},Y^{(n)} \sim \text{Unif}( \partial B_p^n)$ or $ X^{(n)},Y^{(n)} \sim \text{Unif}(B_p^n)$ be independent random vectors. Then,
\begin{align*}
\sqrt{n} \left( n^{1/p - 1/2} \Vert X^{(n)} - Y^{(n)} \Vert _2 -  \sqrt{2} \sqrt{\frac{\Gamma (3/p)}{\Gamma(1/p)}} p^{1/p}  \right) \overset{d}{\longrightarrow} N(0, \sigma_p ^2),
\end{align*}
where $\sigma_p ^2>0$ is given by
\begin{align*}
\sigma_p ^2 = \frac{p^{2/p} \left( M_p(4) + ((M_p(2))^2 \right)}{4 M_p(2)}.
\end{align*}
\end{theorem}

Note that, due to the nature of central limit theorems, the same CLT holds when we consider $X^{(n)}, Y^{(n)}$ to be uniformly distributed on $l_p^n$-balls or on their boundary. Independently, we obtain a similar result for the $n$-dimensional cube and its boundary. Moreover, taking the limit as $p \to \infty$ in Theorem \ref{Thm:CLTforlpballs}, it recovers the following result.

\begin{theorem}
\label{Thm:CLTforCube}
Let $X^{(n)},Y^{(n)} \sim \text{Unif}( \partial B_\infty ^n)$ or $X^{(n)},Y^{(n)} \sim \text{Unif}(B_\infty ^n)$. Then,
\begin{align*}
\sqrt{n} \left( n^{-1/2}  \Vert X^{(n)} - Y^{(n)} \Vert _2 - \sqrt{\frac{2}{3}}  \right) \overset{d}{\longrightarrow} N \left( 0, \frac{7}{30} \right).
\end{align*}
\end{theorem}



\subsection{Large deviation principles}

Large deviation principles describe the asymptotic behavior of the probability of rare events for a sequence of random variables by characterizing their exponential decay rate in terms of the speed and a rate function (see section \ref{Subsection:LargeDeviations} for precise definitions).

We now present the main theorems on the scale of large deviations.

\begin{theorem}
\label{ThmLDPBoundarypgeq2}
Let $p\geq 2$, and let $X^{(n)} , Y^{(n)} \sim \text{Unif} ( \partial B_p^n)$ be independent random vectors. Then, the sequence $\left( \dist (X^{(n)} , Y^{(n)}) \right)_{n \in \N} := \left( n^{1/p - 1/2} \Vert X^{(n)} - Y^{(n)} \Vert _2 \right)_{n \in \N}$ satisfies an LDP with speed $n$ and good rate function
\[   
I_{\dist (X^{(n)} , Y^{(n)})}(z) = 
     \begin{cases}
       \displaystyle{\inf_{ \substack{x \geq 0, y > 0 \\ x^{1/2}y^{-1/p} = z}} \Lambda ^* (x,y)}  &\quad\text{: }z \geq 0\\
       + \infty &\quad\text{: }z < 0, \\ 
     \end{cases}
\]
where $ \Lambda^* $ is the Legendre-Fenchel transform of the function
\begin{align*}
\Lambda (t_1,t_2) & = \log \int _{-\infty} ^{\infty} \int _{-\infty} ^{\infty} e^{t_1 |x-y|^2 + t_2|y|^p    
  - \frac{|x|^p}{p} - \frac{|y|^p}{p} } \frac{1}{4p^{2/p} \left( \Gamma(1 + 1/p) \right)^2} dx dy .
\end{align*}
\end{theorem}

\begin{rmk}
Let us point out that $\dist (\cdot , \cdot)$ denotes a normalised distance rather than the standard Euclidean distance. Nonetheless, they coincide if $p=2$.
\end{rmk}

In contrast to the CLT, large deviation principles exhibit a more sensitive dependence on the underlying distribution. As shown in the following theorems, the LDPs for points uniformly distributed on the boundary and those in $l_p^n$-balls lead to different behaviors, reflecting the focus of large deviations on tail events rather than fluctuations around the mean.

\begin{theorem}
\label{ThmLDPpgeq2}
Let $p\geq 2$, and let $X^{(n)} , Y^{(n)} \sim \text{Unif} ( B_p^n)$ be independent. Then, the sequence $\left( \text{dist}(X^{(n)} , Y^{(n)}) \right)_{n \in \N} := \left( n^{1/p - 1/2} \Vert X^{(n)} - Y^{(n)} \Vert _2 \right)_{n \in \N}$ satisfies an LDP with speed $n$ and good rate function
\[   
I_{\dist(X^{(n)} , Y^{(n)})}(z) = 
     \begin{cases}
       \displaystyle{\inf_{ \substack{z_1 \geq 0, z_2 \geq 0 \\ z=z_1z_2}} I_V (z_1,z_2)}  &\quad\text{: }z \geq 0\\
       + \infty &\quad\text{: } z<0, \\ 
     \end{cases}
\]
where
\begin{align*}
I_V(z_1,z_2) := I_U(z_1) + I_W(z_2), \quad (z_1,z_2)\in \R^2,
\end{align*}
being
\[   
I_U(z) = 
     \begin{cases}
       \log (z) &\quad\text{: }z \in (0, 1]\\
       + \infty &\quad\text{: otherwise } \\ 
     \end{cases}
\]
and
\[   
I_W(z) = 
     \begin{cases}
       \displaystyle{\inf_{ \substack{x \geq 0, y > 0 \\ x^{1/2}y^{-1/p} = z}} \Lambda ^* (x,y)}  &\quad\text{: }z \geq 0\\
       + \infty &\quad\text{: }z < 0, \\ 
     \end{cases}
\]
where $ \Lambda^* $ is the Legendre-Fenchel transform of the function
\begin{align*}
\Lambda (t_1,t_2) & = \log \int _{-\infty} ^{\infty} \int _{-\infty} ^{\infty} e^{t_1 |x-y|^2 + t_2|y|^p    
  - \frac{|x|^p}{p} - \frac{|y|^p}{p} } \frac{1}{4p^{2/p} \left( \Gamma(1 + 1/p) \right)^2} dx dy .
\end{align*}
\end{theorem}


We now turn to the case of the $n$-dimensional cube, where we establish a large deviations principle for points uniformly distributed in the cube.

\begin{theorem}
\label{ThmLDPBoundaryCube}
Let $X^{(n)} , Y^{(n)} \sim \text{Unif} ( B_\infty^n)$ be independent random vectors. Then, the sequence $\left( \text{dist}(X^{(n)} , Y^{(n)}) \right)_{n \in \N} := \left( n^{ - 1/2} \Vert X^{(n)} - Y^{(n)} \Vert _2 \right)_{n \in \N}$ satisfies an LDP with speed $n$ and good rate function
\[   
I_{\text{dist}(X^{(n)} , Y^{(n)})}(z) = 
     \begin{cases}
       \displaystyle{\inf_{ \substack{x \geq 0 \\ z = \sqrt{x}}} \Lambda^* (x)}  &\quad\text{: }z \geq 0\\
       + \infty &\quad\text{: }z<0 , \\ 
     \end{cases}
\]
where $ \Lambda^* $ is the Legendre-Fenchel transform of the function
\begin{align*}
\Lambda (t) = \log \left[ 2 \int _{0} ^2 e^{t x^2 } \left( 1- \frac{x}{2} \right) dx \right] .
\end{align*}
\end{theorem}


The rest of this paper is structured as follows. Some notation and results needed for the main results are collected in Section \ref{section2}. In Section \ref{section3}, we introduce the problem considered and prove the spherical case following the idea of \cite{H} in a more direct way.  Section \ref{section4} contains the proofs of the central limit theorems (Theorems \ref{Thm:CLTforlpballs} and \ref{Thm:CLTforCube}) and Section \ref{section5} contains the proofs of the large deviation principles (Theorems \ref{ThmLDPBoundarypgeq2}, \ref{ThmLDPpgeq2} and \ref{ThmLDPBoundaryCube}).

 \section{Notation and preliminaries}
 \label{section2}

\subsection{General notation}

Let $A \subset \mathbb{X}$ be a subset of some topological space $\mathbb{X}$, we denote its interior as $A^\circ$ and its closure as $\overline{A}$. Given a Borel probability measure $\tau$ on $\R^n$, we indicate by $X \sim \tau$ that the random vector $X$ has distribution $\tau$. In particular, we denote as $X \sim \mathcal{N}( \mu , \sigma ^2 )$ that the random variable $X$ has a Gaussian distribution with mean $\mu \in \R$ and variance $\sigma ^2 > 0$. In dimension $n \geq 2$, we write $\mathcal{N} ( \mu , \Sigma)$ to indicate the multivariate Gaussian distribution with mean $\mu \in \R^n$ and covariance matrix $\Sigma \in \R^{n \times n} $. Given $(X_n)_{n \in \N}$ a sequence of random vectors in $\R^n$, we write $ \overset{d}{\to}$, $ \overset{p}{\to}$ and $ \overset{a.s.}{\to}$ to indicate convergence in distribution, probability and almost sure, respectively. Moreover, given two random elements $X_1$ and $X_2$, we indicate equality in distribution as $X_1 \overset{d}{=} X_2$. Given two sequences of random vectors $(X_n)_{n \in \N}$, $(Y_n)_{n \in \N}$, we denote asymptotic equivalence in distribution, i.e., $\Vert X_n - Y_n \Vert _2 \overset{p}{\to} 0$, as $X_n \overset{d}{\sim} X_n$. We write $\E$ and $\var$ for the expected value and variance, respectively.


\subsection{Probabilistic aspects of $l_p^n$-balls}
\label{Section2.2}

Recall that for $p \in [1, \infty ]$ and $x = (x_1, ... , x_n) \in \R^n$, $n \in \N$, we write $\Vert x \Vert _p$ for the $p$-norm of $x$. We denote the unit ball in $\R^n$ with respect to the $\Vert \cdot \Vert _p$-norm by $B_p^n$, that is 
\begin{align*}
B_p^n := \lbrace x \in \R^n \; \Vert x \Vert _p \leq 1 \rbrace,
\end{align*}
and we write $\partial B_p^n$ for its boundary. Additionally, we will write $S^{n-1}:= \partial B_2^n$ to denote the Euclidean sphere.

Let $\sigma _p ^n$ be the surface measure (Hausdorff measure) on $\partial B_p^n$, $p \geq 1$, and denote by $\mu_p ^n$ the cone probability measure on $\partial B_p^n$, defined by
\begin{align*}
\mu_p ^n (A) = \frac{1}{\vol (A)} \vol \left( \lbrace ta \in \R^n ; a \in A, 0\leq t \leq 1 \rbrace \right), \; \; \; \;  A \subset \partial B_p^n,
\end{align*}
where $\vol$ denotes the Lebesgue measure.

Recall that a random variable is said to be $p-$generalized Gaussian ($1\leq p < \infty$) if its density function $f_p$ is given by
\begin{align*}
f_p(x) = C_p^{-1} e^{- |x|^p / p}, \; \; \; x \in \R,
\end{align*}
with respect to the Lebesgue measure on $\R$, where the normalization constant $C_p$ is given by $C_p:= 2 p^{1/p} \Gamma \left(1 + \frac{1}{p} \right)$. For the case $p=\infty$, we define the $\infty$-generalized Gaussian as a random variable uniformly distributed on $[-1,1]$. Notice that, if $g$ is a $p$-generalized Gaussian for some $p \in [1, \infty)$, then for every $\alpha \geq 0$, $\E | g |^\alpha = M_p(\alpha )$.

We recall from \cite{SZ} the following probabilistic representation for a uniformly distributed random vector in $B_p^n$ and in $\partial B_p^n$.

\begin{lemma}[Schechtman and Zinn]
\label{Lem:SchAndZinn}
Fix $1\leq p < \infty$. Let $Y = (Y_1, ... , Y_n)$ be a random vector with independent $p$-generalized Gaussian coordinates $Y_1,...,Y_n$, and let $U$ be a random variable uniformly distributed on $[0,1]$, independent of $Y$. Then,

\begin{itemize}

\item The random vector $\displaystyle{\frac{Y}{\Vert Y \Vert _p}}$ and the random varianble $\Vert Y \Vert _p$ are independent.

\item $\displaystyle{\frac{Y}{\Vert Y \Vert _p}}$ is uniformly distributed on $\partial B_p^n$ according to the cone measure $\mu _p^n$.

\item $\displaystyle{U^{1/n} \frac{Y}{\Vert Y \Vert _p}}$ is uniformly distributed in $B_p^n$.

\end{itemize}
\end{lemma}

\subsection{Delta Method} 
We will use the following technical lemma, which characterizes the asymptotic distribution of a differentiable function of a random vector which is asymptotically Gaussian (see \cite{O}).

\begin{lemma}[Delta method]
\label{DeltaMethod}
Let $(X_n)_{n\in\N}$ be a sequence of $k-$dimensional random vectors such that $\sqrt{n}(X_n - \mu)$ converges in distribution to a centered Gaussian random vector with covariance matrix $\Sigma$, and let $g:\mathbb{R}^k \to \R^d$ be continuously differentiable at $\mu$ with Jacobian matrix $J_g$. Then,
\begin{align*}
\sqrt{n} (g(X_n) - g(\mu) ) \overset{d}{\longrightarrow} N\left( 0 , J_g \Sigma J_g^T \right)
\end{align*}
provided $J_g \Sigma J_g^T$ is positive definite.
\end{lemma}


\subsection{Large deviations}
\label{Subsection:LargeDeviations}
Let $X := (X_n)_{n \in \N}$ be a sequence of random elements taking values in some Hausdorff topological space $\mathbb{X}$. Then, we say that $X$ satisfies a large deviation principle (LDP) with speed $s(n)$ and rate function $I_{X}$ if $s: \N \to (0, \infty)$, $I_{X}: \mathbb{X} \to [0, \infty]$ is lower semi-continuous, and if
\begin{align*}
- \inf _{x \in A^\circ } I_X (x) & \leq \liminf _{n \to \infty} \frac{1}{s(n)} \log \mathbb{P} (X_n \in A ) \leq \limsup _{n \to \infty} \frac{1}{s(n)} \log \mathbb{P} (X_n \in A ) \leq - \inf _{x \in \overline{ A} } I_X (x)
\end{align*}
for all Borel sets $A \subset \mathbb{X}$. 

The rate function $I_X$ is said to be a good rate function if it has compact level sets $\lbrace x \in \mathbb{X} ; I_X (x) \leq \alpha \rbrace$, for every $\alpha \in [0, \infty)$.

Let $\mathbb{X} = \R^d$ for some $d \in \N$ and let $\Lambda : \R^d \to \R \cup \lbrace +\infty \rbrace$. Then, the Legendre-Fenchel transform of $\Lambda$ is denoted as $\Lambda ^* $ and defined as
\begin{align*}
\Lambda^* (x) := \sup _{u \in \R^d} \left[ \langle x , u \rangle - \Lambda (u) \right], \; \; x \in \R^d,
\end{align*}
where $\langle \cdot , \cdot \rangle$ denotes the standard scalar product on $\R^d$. Moreover, we define the (effective) domain of $\Lambda$ to be the set $D_\Lambda := \lbrace u \in \R^d ; \Lambda (u) < \infty \rbrace$.

We now present some fundamental results from large deviations theory. First, we recall what is known as Cram\'er's theorem, which provides an LDP for sequences of independent and identically distributed random vectors (Theorem 2.2.3. in \cite{DZ}).
\begin{lemma}[Cramer's theorem]
Let $X_1,X_2,...$ be independent and identically distributed random vectors taking values in $\R^d$. Assume that the origin is an interior point of $D_\Lambda$, where $\Lambda (u) = \log \E e^{\langle u , X \rangle }$. Then, the sequence of partial sums $\left( \frac{1}{n} \sum _{i=1} ^n X_i \right)_{n \in \mathbb{N}}$, satisfies an LDP on $\R^n$ with speed $n$ and good rate function $\Lambda ^* $.
\end{lemma}

Next result, which will be used in the following sections, states that if $(X_n)_{n \in \N}$ is a sequence of random variables that satisfies an LDP with speed $s(n)$ and rate function $I_X$ and $(Y_n)_{n \in \N}$ is a sequence of random variables that is "close" to $(X_n)_{n \in \N}$ in some sense, then the sequence $(Y_n)_{n \in \N}$ will also satisfy an LDP under some conditions (Theorem 4.2.13. in \cite{DZ}).
\begin{lemma}[Exponential equivalence]
\label{LemExponentialEquivalence}
Let $X = (X_n)_{n \in \N}$ and $Y = (Y_n)_{n \in \N}$ be two sequences of random variables and assume that $X$ satisfies an LDP with speed $s(n)$ and rate function $I_X$. Further, assume that $X$ and $Y$ are exponentially equivalent, i.e., 
\begin{align*}
\limsup _{n \to \infty} \frac{1}{s(n)} \log \mathbb{P} ( | X_n - Y_n | > \delta ) = - \infty
\end{align*}
for any $\delta >0$. Then, $Y$ satisfies an LDP with the same speed and the same rate function as $X$.
\end{lemma}

Lastly, the contraction principle allows to transform a given LDP to another by a continuous function (Theorem 4.2.1. in \cite{DZ}). 
\begin{lemma}[Contraction principle]
\label{LemContractionPrinciple}
Let $\mathbb{X}$ and $\mathbb{Y}$ be two Hausdorff topological spaces and $F: \mathbb{X} \to \mathbb{Y}$ be a continuous function. Further, let $X = (X_n)_{n \in \N}$ be a sequence of $\mathbb{X}$-valued random elements that satisfies an LDP with speed $s(n)$ and good rate function $I_X$. Then, the sequence $Y:=  ( F(X_n) )_{n \in \N}$ satisfies an LDP on $\mathbb{Y}$ with the same speed and with good rate function $I_Y = I_X \circ F^{-1}$, i.e., $I_Y(y) := \inf \lbrace I_X(x) ; F(x) = y \rbrace$, $y \in \mathbb{Y}$ with the convention that $I_Y(y) = +\infty$ if $F^{-1} ( \lbrace y \rbrace ) = \emptyset$. 
\end{lemma}



%\begin{lemma}[Cramer's theorem for stretched exponentials]
%Let $X,X_1,X_2,... $ be non-negative, independent and identically distributed random variables. Assume that there are constants $r \in (0,1)$, $t_0 >0$ and slowly varying functions $c_1,c_2,b : (0,\infty ) \to (0, \infty)$ such that
%\begin{align*}
%c_1(z) e^{-b(z) z^{p/q}} \leq \mathbb{P} \left( |g_1 | \geq z \right) \leq c_2(z) e^{-b(z) z^{p/q}}
%\end{align*}
%\end{lemma}










\section{Previous results and comparison}
\label{section3}

In this section, we will provide a short and compact proof of Hammersley's CLT in the case of the $n$-dimensional Euclidean sphere (see \cite{H}).

Let $X^{(n)} , Y^{(n)}$ be independent random vectors uniformly distributed on the Euclidean $n$-dimensional sphere, $S^{n-1}$. In \cite{H}, the author approached the problem by studying the distribution function of the distance between $X^{(n)}$ and $Y^{(n)}$, which was shown to be close to an incomplete beta-function ratio, and therefore, tends asymptotically to a Gaussian, as the dimension grows.

Following \cite{H}, in \cite{L}, the author proved the result by extending it to spherical distributions, whose density function depends only on the modulus of the random vector.

The underlying idea of the proofs relies on the rotational invariance of the sphere, which implies that $\dist(X^{(n)},Y^{(n)})$ depends only on a single random variable. To see this, fix a point $x_0 \in S^{n-1}$ on the sphere. For any $x\in S^{n-1}$,
$$ \int _{S^{n-1}} \mathbf{1}_{ \lbrace \dist(x,y) \rbrace }  d \sigma(y) =  \int _{S^{n-1}} \mathbf{1}_{ \lbrace \dist(x_0,y) \rbrace } d\sigma(y), $$
where $\sigma$ is the normalized cone measure on $S^{n-1}$, i.e., $\sigma (S^{n-1}) = 1$. 
Then, 
\begin{align*}
\mathbb{P} (\dist (X^{(n)}, Y^{(n)}) \leq t) & = \int_{S^{n-1}} \mathbb{P} (\dist(x, Y^{(n)}) \leq t \mid X^{(n)} = x)  \, d\sigma(x)
 \\
 & =  \int_{S^{n-1}} \int_{S^{n-1}} \mathbf{1}_{\{\dist(x, y) \leq t\}} \, d \sigma(y) \, d\sigma(x)
 \\
 & =  \int_{S^{n-1}}  \mathbf{1}_{\{\dist(x_0, y) \leq t\}} \, d\sigma(y) \int_{S^{n-1}}  d\sigma(x)
 \\
 & =  \int_{S^{n-1}} \mathbf{1}_{\{\dist(x_0, y) \leq t\}} \, d\sigma(y)
 \\
 & = \mathbb{P} (\dist(x_0, Y^{(n)}) \leq t).
\end{align*}



Therefore, the problem is reduced to studying the asymptotics of the random variable $\dist(x_0, Y^{(n)})$. For the general case, i.e., when $X^{(n)} , Y^{(n)}$ are uniformly distributed on $B_p^ n$ or $\partial B_p^n$, for $p\in [1, \infty]$ and $n \in \N$, we develop suitable techniques in sections \ref{section4} and \ref{section5}.

Let us prove the following theorem.
\begin{theorem}
\label{Thm:CLThypersphere}
Let $X^{(n)},Y^{(n)} \sim \text{Unif} (S^{n-1})$ be independent random vectors. Then,
\begin{align*}
  \frac{ \dist(X^{(n)},Y^{(n)}) - \mathbb{E} [\dist(X^{(n)},Y^{(n)})]}{ \sqrt{ Var[\dist(X^{(n)},Y^{(n)})]} } \overset{d}{\to} N(0,1) .
\end{align*}
In particular, 
\begin{align*}
\sqrt{2n} \left( \dist(X^{(n)},Y^{(n)}) - \sqrt{2} \right) \overset{d}{\to}  N(0,1) .
\end{align*}
\end{theorem}

To prove this result, let us first compute the distribution function of the distance.

\begin{lemma}
\label{lem:DistribDistSphere}
Let $X^{(n)},Y^{(n)}$ be independent random vectors uniformly distributed on $S^{n-1}$, for each $n \in \N$. Then, for every $t \in [0,2]$
\begin{align*}
\mathbb{P} (\dist(X^{(n)},Y^{(n)}) \leq t) & = \frac{\Gamma ( \frac{n}{2} ) }{ \sqrt{\pi} \Gamma ( \frac{n-1}{2} )} \int_{1-\frac{t^2}{2}}^{1} (1- x^2)^{(n-3)/2} dx 
 \\
  & = \frac{\Gamma ( \frac{n}{2} ) }{ \sqrt{\pi} \Gamma ( \frac{n-1}{2} )} \int_0^{s}  (\sin (\theta))^{n-2} d\theta ,
\end{align*}
where $s = arc cos (1-t^2/2)$.
\end{lemma}
\begin{proof}

Without loss of generality, we fix a point on the sphere $x_0 = e_1$. Let $\theta \in [0,2\pi ]$ be the angle between the line joining $x_0$ and the origin, with the line joining the origin and a point $y_\theta$ on the sphere.

First, note that the lenght of the chord joining $x_0$ and $y_\theta$ is $2 \sin (\theta / 2)$, and so it follows that $\cos( \theta) = 1- \dist(x_0 , y _\theta)^2 / 2 $.

Given an angle $\theta \in [0, 2\pi ]$, the $(n-2)$-dimensional Hausdorff measure of the points at distance $2 \sin (\theta / 2)$ of $x_0$ is
$$
\frac{2 \pi^{\frac{(n-1)}{2}}}{\Gamma ( \frac{n-1}{2} )} (\sin (\theta))^{n-2} .
$$
Notice that, given $t \in [0,2]$, the event $\lbrace \dist (x_0,Y^{(n)}) \leq t \rbrace$ has the same pobability as the event $\lbrace 2\sin ( \theta / 2 ) \leq t \rbrace$, where $\theta$ is again the angle between the line joining $x_0$ and the origin, with the line joining the origin and $Y^{(n)}$.

Therefore,
\begin{align*}
\mathbb{P} ([\dist(X^{(n)},Y^{(n)})] \leq t) & = \int_0^{s} \frac{\frac{2 \pi^{\frac{n-1}{2}}}{\Gamma ( \frac{n-1}{2} )} (\sin (\theta))^{n-2}}{\frac{2 \pi^{\frac{n}{2}}}{\Gamma ( \frac{n}{2} )}} d\theta 
 \\
 & = \frac{\Gamma ( \frac{n}{2} ) }{ \sqrt{\pi} \Gamma ( \frac{n-1}{2} )} \int_0^{s}  (\sin (\theta))^{n-2} d\theta ,
\end{align*}
where $t = 2 \sin (s / 2)$ or $t = \sqrt{2} \sqrt{1- \cos (s)}$, i.e., $s = 2 \arcsin (t / 2)$ or $s = \arccos (1-t^2/2)$. Thus, we can also write is as
\begin{align*}
\mathbb{P} ([\dist(X^{(n)},Y^{(n)})] \leq t) & = \frac{\Gamma ( \frac{n}{2} ) }{ \sqrt{\pi} \Gamma ( \frac{n-1}{2} )} \int_{1-\frac{t^2}{2}}^{1} (1- x^2)^{(n-3)/2} dx.
\end{align*}
\end{proof}

Now, we are ready to compute the expected value of the distance.

\begin{lemma}
\label{Lem:ExpectedValueHypersphere}
Let $X^{(n)},Y^{(n)} \sim \text{Unif}( S^{n-1})$ independent random vectors. Then, 
$$ \mathbb{E} [ \dist(X^{(n)},Y^{(n)})] = \frac{2^{n-1} \left( \Gamma(n/2) \right)^2}{\sqrt{\pi}\Gamma(n-1/2)}  \underset{n \to \infty}{ \longrightarrow} \sqrt{2}. $$
\end{lemma}
\begin{proof}
By differentiating the distribution function of $\dist ( X^{(n)} , Y^{(n)})$, we obtain its density function given by
\begin{align*}
f(t) = \frac{\Gamma ( \frac{n}{2} ) }{ \sqrt{\pi} \Gamma ( \frac{n-1}{2} )} \left( \sin \left( 2 \arcsin \left( \frac{t}{2} \right) \right) \right)^{n-2} \frac{1}{\sqrt{1-\frac{t^2}{4}}} .
\end{align*}
Applying the change of variables $\theta = 2 \arcsin ( t/2 )$, we obtain that
\begin{align*}
\mathbb{E}[\dist(x_0 , Y^{(n)})] & = \frac{\Gamma ( \frac{n}{2} ) }{ \sqrt{\pi} \Gamma ( \frac{n-1}{2} )} \int_0^{\pi} 2\sin{(\theta/2)} (\sin (\theta))^{n-2} d\theta  .
\end{align*}
Using the identity $\sin (x) = 2 \sin (x/2) \cos (x/2)$ and the change of variables $x=\theta/2$,
\begin{align*}
\int_0^{\pi} 2\sin{(\theta/2)} (\sin (\theta))^{n-2} d\theta & = 
\int_0^{\pi} 2\sin{(\theta/2)} (2 \sin (\theta/2) \cos (\theta/2))^{n-2} d\theta
 \\
 & = \int_0^{\pi} ( 2\sin{(\theta/2)} ) ^{n-1} ( \cos ( \theta /2))^{n-2} d\theta
 \\
 & = 2^{n} \int_0^{\frac{\pi}{2}} ( \sin ( x ) ) ^{n-1} ( \cos (x))^{n-2} dx
  \\
  & = 2^{n-1} B \left( \frac{n}{2} , \frac{n-1}{2}  \right),
\end{align*}
where $B$ is the Beta function. Thus, we conclude that 
\begin{align*}
\mathbb{E}[\dist(x_0 , Y^{(n)})] & = \frac{\Gamma ( \frac{n}{2} ) }{ \sqrt{\pi} \Gamma ( \frac{n-1}{2} )} 2^{n-1} B \left( \frac{n}{2} , \frac{n-1}{2}  \right) 
 \\
 & = \frac{\Gamma ( \frac{n}{2} ) }{ \sqrt{\pi} \Gamma ( \frac{n-1}{2} )} 2^{n-1} 
 \frac{\Gamma ( \frac{n}{2} ) \Gamma ( \frac{n-1}{2} ) }{ \Gamma ( \frac{2n - 1}{2} )}
 \\
 & = \frac{2^{n-1} \left( \Gamma(n/2) \right)^2}{\sqrt{\pi}\Gamma(n-1/2)}.
\end{align*}
Furthermore, using the second order Stirling's formula,
\begin{align*}
\Gamma (1+n) = \sqrt{2 \pi n } \left( \frac{n}{e} \right)^n \left( 1 + \frac{1}{12n} + O\left( \frac{1}{n^2} \right) \right),
\end{align*}
and applying it to $\Gamma(n/2)$ and $\Gamma ( n- 1/2 )$, a straightforward computation yields  $$\mathbb{E}[\dist(x_0 , Y^{(n)})] = \sqrt{2} - \frac{1}{4 \sqrt{2}n} + o\left( 1/n \right) .$$
%& = \frac{2^n [ \Gamma((n+1)/2)]^2}{\sqrt{\pi}\Gamma(n+1/2)}
 %\\
 %& =  \frac{2  \Gamma(n) \Gamma((n+1)/2)}{ \Gamma( n/2) \Gamma(n+1/2)}.
\end{proof}

In a similar way, we can compute the variance of the distance and its asymptotics.
\begin{lemma}
\label{Lem:VarianceHypersphere}
Let $X^{(n)},Y^{(n)} \sim \text{Unif}( S^{n-1})$ independent random vectors. Then, 
$$ Var[\dist(X^{(n)},Y^{(n)})] = 2 - \left( \frac{2^{n-1} \left( \Gamma(n/2) \right)^2}{\sqrt{\pi}\Gamma(n-1/2)} \right)^{2} , $$
and 
$$  \lim _{n \to \infty} 2 n  Var [\dist(X^{(n)} , Y^{(n)})] = 1. $$
\end{lemma}
\begin{proof}
Let us compute
$$
Var [\dist(x_0 , Y^{(n)}) ] = 
\mathbb{E}[(\dist(x_0 , Y^{(n)}))^2] - \left( \mathbb{E}[\dist(x_0 , Y^{(n)})] \right)^2.
$$
Since for every $\theta \in [ 0, 2\pi ]$ $$
\sin (\theta / 2 ) = \sqrt{\frac{1-\cos \theta}{2}},
$$
we have that
\begin{align*}
\mathbb{E} \left[ \left( \dist(x_0 , Y^{(n)})\right)^2 \right] & = \frac{\Gamma ( \frac{n}{2} ) }{ \sqrt{\pi} \Gamma ( \frac{n-1}{2} )} \int_0^{\pi} (2\sin((\theta/2))^2 (\sin (\theta))^{n-2} d\theta 
 \\
 & = \frac{2 \Gamma (  \frac{n}{2} ) }{ \sqrt{\pi} \Gamma ( \frac{n-1}{2} )} \left[ \int_0^{\pi} (\sin (\theta))^{n-2}  d\theta
 - \int_0^{\pi} \cos (\theta) (\sin (\theta))^{n-2} d\theta  \right]
\end{align*}
On the one hand, using the Beta function,
$$
\int_0^{\pi} (\sin (\theta))^{n-2}  d\theta =
2 \int_0^{\frac{\pi}{2}} (\sin (\theta))^{n-2}  d\theta 
  =  B \left( \frac{n-1}{2} , \frac{1}{2}  \right) = \frac{\sqrt{\pi} \Gamma \left( \frac{n-1}{2} \right)}{ \Gamma(n/2) }.
$$
On the other hand,
%\begin{align*}
%\int_0^{\pi} \cos (\theta) (\sin (\theta))^{n-2} d\theta  = \frac{2}{n-1}
%\end{align*}
\begin{align*}
\int_0^{\pi} \cos (\theta) (\sin (\theta))^{n-2} d\theta  = 0.
\end{align*}
Therefore,
\begin{align*}
Var [\dist(x_0 , Y^{(n)}) ] & = \frac{2 \Gamma ( \frac{n}{2} ) }{ \sqrt{\pi} \Gamma ( \frac{n-1}{2} )} \left[ \frac{\sqrt{\pi} \Gamma ( \frac{n-1}{2} )}{ \Gamma(n/2) } \right] - \left( \frac{2^{n-1} \left( \Gamma(n/2) \right)^2}{\sqrt{\pi}\Gamma(n-1/2)} \right)^2
 \\
 & = 2 - \left( \frac{2^{n-1} \left( \Gamma(n/2) \right)^2}{\sqrt{\pi}\Gamma(n-1/2)} \right)^{2}.
\end{align*}
Furthermore, by Lemma \ref{Lem:ExpectedValueHypersphere} taking the limit $n \to \infty$, we obtain that
$$ \lim _{n \to \infty} Var [\dist(X^{(n)} , Y^{(n)})] = 0  $$
and
$$ \lim _{n \to \infty} 2n  Var [\dist(X^{(n)} , Y^{(n)})]= 1 .  $$
\end{proof}

Before proving Theorem \ref{Thm:CLThypersphere}, we need the following technical result (see \cite{S}, Theorem).

\begin{theorem}
\label{UniformIntegrabTheorem}
Let $( p_n )_{n \in \N}$ be a sequence of densities on $\R$ such that, 
\begin{align*}
\lim _{n \to \infty} p_n (x) = p (x)
\end{align*}
for almost every $x \in \R$. If $p$ is a density, then
\begin{align*}
\lim _{n \to \infty} 	\int_ S p_n(x) dx = 	\int_ S p(x) dx 
\end{align*}
uniformly for all Borel sets $S$ in $\R$.
\end{theorem}

Now, we are ready to prove Theorem \ref{Thm:CLThypersphere}.
\begin{proof}[Proof of Theorem \ref{Thm:CLThypersphere}]
We write $\mu _n := \mathbb{E} [\dist(X^{(n)},Y^{(n)})]$ and $\sigma^2 _n := Var[\dist(X^{(n)},Y^{(n)})]$, whose values are provided by Lemmas \ref{Lem:ExpectedValueHypersphere} and \ref{Lem:VarianceHypersphere}. Then, note that
\begin{align*}
\frac{ \dist(X^{(n)},Y^{(n)}) -\mu _n}{ \sigma _n } = \frac{\dist(X^{(n)},Y^{(n)}) - \sqrt{2}}{\frac{1}{\sqrt{2n}}} \frac{1}{\sqrt{2n} \sigma _n} + \frac{\sqrt{2} - \mu_n}{\sigma _n}.
\end{align*}
By Slutsky's Theorem, to prove Theorem \ref{Thm:CLThypersphere}, it is enough to prove that
\begin{align*}
\frac{\dist(X^{(n)},Y^{(n)}) - \sqrt{2}}{\frac{1}{\sqrt{2n}}} \overset{d}{ \to } N(0,1),
\end{align*}
since by Lemmas \ref{Lem:ExpectedValueHypersphere} and \ref{Lem:VarianceHypersphere},
we have that $$ \left| \frac{\sqrt{2} - \mu _n}{\sigma_n} \right|  = \left| \frac{1}{4\sqrt{2}n \sigma _n} + \frac{1}{\sigma _n} o \left( \frac{1}{n} \right) \right| $$
and 
$$\lim _{n \to \infty} \sigma _n \sqrt{2n} = 1.$$

Note that the random variable $\dist(X^{(n)},Y^{(n)})$ takes values on the interval $[0,2]$. Thus, the normalized random variable $\sqrt{2n} \left( \dist(X^{(n)},Y^{(n)}) - \sqrt{2} \right) $ takes values of the interval $[-2\sqrt{n},2\sqrt{2n} -2\sqrt{n}  ]$.

By Lemma \ref{lem:DistribDistSphere} we have that for any $t \in [-2\sqrt{n},2\sqrt{2n} -2\sqrt{n}  ]$,
\begin{align*}
\mathbb{P} \left(\frac{\dist(X^{(n)},Y^{(n)}) - \sqrt{2}}{\frac{1}{\sqrt{2n}}} \leq t\right) & =
 \frac{\Gamma ( \frac{n}{2} ) }{ \sqrt{\pi} \Gamma ( \frac{n-1}{2} )} 
 \int_{ -\frac{t^2}{4n} - \frac{t}{\sqrt{n}}      }^{1} (1- x^2)^{(n-3)/2} dx .
\end{align*}
Let $t \in [-2\sqrt{n},2\sqrt{2n} -2\sqrt{n} ]$. We use the change of variables $\sqrt{n} x = y$, and obtain
\begin{align*}
\mathbb{P} \left(\frac{\dist(X^{(n)},Y^{(n)}) - \sqrt{2}}{\frac{1}{\sqrt{2n}}} \leq t\right) & =
 \frac{\Gamma ( \frac{n}{2} ) }{ \sqrt{\pi} \Gamma ( \frac{n-1}{2} ) \sqrt{n}} 
 \int_{ -\frac{t^2}{4\sqrt{n}} - t      }^{\sqrt{n}} \left(1- \frac{y^2}{n} \right)^{(n-3)/2} dy.
\end{align*}
Let $t \in \mathbb{R}$. Let $n_t := \inf \lbrace n  ; t \in [-2\sqrt{n},2\sqrt{2n} -2\sqrt{n}  ] \rbrace$. Then, note that for any $n \geq n_t$, since the integrand is non-negative and bounded above by $1$, we have 
\begin{align*}
& \mathbb{P} \left(\frac{\dist(X^{(n)},Y^{(n)}) - \sqrt{2}}{\frac{1}{\sqrt{2n}}} \leq t\right) \cr
& =  \frac{\Gamma ( \frac{n}{2} ) }{ \sqrt{\pi} \Gamma ( \frac{n-1}{2} ) \sqrt{n}} 
 \int_{ -\frac{t^2}{4\sqrt{n}} - t}^{\sqrt{n}} \left(1- \frac{y^2}{n}\right)^{(n-3)/2} \mathds{1}_{[-\sqrt{n} , \sqrt{n}]} (y) dy   \cr
& = \frac{\Gamma ( \frac{n}{2} ) }{ \sqrt{\pi} \Gamma ( \frac{n-1}{2} ) \sqrt{n}}  \Bigg(  \int_{ -t}^{\sqrt{n}} \left(1- \frac{y^2}{n}\right)^{(n-3)/2} \mathds{1}_{[-\sqrt{n} , \sqrt{n}]} (y) dy   + \int_{-\frac{t^2}{4\sqrt{n}} - t }^{-t} \left(1- \frac{y^2}{n}\right)^{(n-3)/2} \mathds{1}_{[-\sqrt{n} , \sqrt{n}]} (y)  dy  \Bigg) \cr
& = \frac{\Gamma ( \frac{n}{2} ) }{ \sqrt{\pi} \Gamma ( \frac{n-1}{2} ) \sqrt{n}}  \Bigg(  \int_{-t}^{\infty} \left(1- \frac{y^2}{n}\right)^{(n-3)/2} \mathds{1}_{[-\sqrt{n} , \sqrt{n}]} (y) dy + \int_{-\frac{t^2}{4\sqrt{n}} - t }^{-t} \left(1- \frac{y^2}{n}\right)^{(n-3)/2} \mathds{1}_{[-\sqrt{n} , \sqrt{n}]} (y) dy \Bigg) \cr
& \leq  \frac{\Gamma ( \frac{n}{2} ) }{ \sqrt{\pi} \Gamma ( \frac{n-1}{2} ) \sqrt{n}}  \Bigg(  \int_{-t}^{\infty} \left(1- \frac{y^2}{n}\right)^{(n-3)/2} \mathds{1}_{[-\sqrt{n} , \sqrt{n}]} (y) dy + \frac{t^2}{4\sqrt{n}} \Bigg).
\end{align*}
We notice that
\begin{align*}
f_n (y) : = \frac{\Gamma ( \frac{n}{2} ) }{ \sqrt{\pi} \Gamma ( \frac{n-1}{2} ) \sqrt{n}}
\left(1- \frac{y^2}{n}\right)^{(n-3)/2} \mathds{1}_{[-\sqrt{n} , \sqrt{n}]} (y)
\end{align*}
is a probability density for every $n \in \N$, and that $f_n \to \frac{1}{\sqrt{2 \pi}} e^{-y^2 / 2}$ pointwise, since taking limits $n \to \infty$, for every $y \in \mathbb{R}$,
$$ \displaystyle{\lim _{n \to \infty} \frac{\Gamma ( \frac{n}{2} ) }{ \sqrt{\pi} \Gamma ( \frac{n-1}{2} ) \sqrt{n}} = \frac{1}{\sqrt{2\pi}}}$$
and
$$ \displaystyle{\lim _{n \to \infty} \left(1- \frac{y^2}{n}\right)^{(n-3)/2} = e^{\frac{-y^2}{2}}} .$$
Then, by Theorem \ref{UniformIntegrabTheorem}, we have that for any $t \in \R$,

\begin{align*}
& \lim _{n \to \infty} \mathbb{P} \left(\frac{\dist(X^{(n)},Y^{(n)}) - \sqrt{2}}{\frac{1}{\sqrt{2n}}} \leq t\right) \cr
& \leq \lim _{n \to \infty} \frac{\Gamma ( \frac{n}{2} ) }{ \sqrt{\pi} \Gamma ( \frac{n-1}{2} ) \sqrt{n}}  \Bigg(  \int_{-t}^{\infty} \left(1- \frac{y^2}{n}\right)^{(n-3)/2} \mathds{1}_{[-\sqrt{n} , \sqrt{n}]} (y) dy + \frac{t^2}{4\sqrt{n}} \Bigg) \cr
& = \lim _{n \to \infty} \frac{\Gamma ( \frac{n}{2} ) }{ \sqrt{\pi} \Gamma ( \frac{n-1}{2} ) \sqrt{n}}   \int_{-t}^{\infty} \lim _{n \to \infty} \left(1- \frac{y^2}{n}\right)^{(n-3)/2} \mathds{1}_{[-\sqrt{n} , \sqrt{n}]} (y) dy   \cr
& = \frac{1}{\sqrt{2 \pi}} 
 \int_{ - t }^{\infty} e^{-\frac{y^2}{2}} dy.
\end{align*}
For the lower bound, note that if $n \geq  n_t$ 
\begin{align*}
& \mathbb{P} \left(\frac{\dist(X^{(n)},Y^{(n)}) - \sqrt{2}}{\frac{1}{\sqrt{2n}}} \leq t\right) \cr
& =  \frac{\Gamma ( \frac{n}{2} ) }{ \sqrt{\pi} \Gamma ( \frac{n-1}{2} ) \sqrt{n}} 
 \int_{ -\frac{t^2}{4\sqrt{n}} - t}^{\sqrt{n}} \left(1- \frac{y^2}{n}\right)^{(n-3)/2} dy   \cr
& = \frac{\Gamma ( \frac{n}{2} ) }{ \sqrt{\pi} \Gamma ( \frac{n-1}{2} ) \sqrt{n}}  \Bigg(  \int_{ -t}^{\sqrt{n}} \left(1- \frac{y^2}{n}\right)^{(n-3)/2} dy   + \int_{-\frac{t^2}{4\sqrt{n}} - t }^{-t} \left(1- \frac{y^2}{n}\right)^{(n-3)/2} dy \Bigg) \cr
& = \frac{\Gamma ( \frac{n}{2} ) }{ \sqrt{\pi} \Gamma ( \frac{n-1}{2} ) \sqrt{n}}  \Bigg(  \int_{-t}^{\infty} \left(1- \frac{y^2}{n}\right)^{(n-3)/2} \mathds{1}_{[-\sqrt{n} , \sqrt{n}]} (y) dy + \int_{-\frac{t^2}{4\sqrt{n}} - t }^{-t} \left(1- \frac{y^2}{n}\right)^{(n-3)/2} dy \Bigg) \cr
& \geq  \frac{\Gamma ( \frac{n}{2} ) }{ \sqrt{\pi} \Gamma ( \frac{n-1}{2} ) \sqrt{n}}   \int_{-t}^{\infty} \left(1- \frac{y^2}{n}\right)^{(n-3)/2} \mathds{1}_{[-\sqrt{n} , \sqrt{n}]} (y) dy .
\end{align*}
Since we already proved that the latter expression converges to $\int_{ - t }^{\infty} e^{-\frac{y^2}{2}} dy$, we obtain that
\begin{align*}
& \lim _{n \to \infty} \mathbb{P} \left(\frac{\dist(X^{(n)},Y^{(n)}) - \sqrt{2}}{\frac{1}{\sqrt{2n}}} \leq t\right) \cr
& \geq  \lim _{n \to \infty} \frac{\Gamma ( \frac{n}{2} ) }{ \sqrt{\pi} \Gamma ( \frac{n-1}{2} ) \sqrt{n}}  \int_{-t}^{\infty} \left( 1- \frac{y^2}{n} \right)^{(n-3)/2} \mathds{1}_{[-\sqrt{n} , \sqrt{n}]} (y) dy \cr
& =  \frac{1}{\sqrt{2 \pi}} 
 \int_{ - t }^{\infty} e^{-\frac{y^2}{2}} dy .
\end{align*}
Thus, for every $t\in \R$,
\begin{align*}
\lim _{n \to\infty } \mathbb{P} \left(\frac{\dist(X^{(n)},Y^{(n)}) - \sqrt{2}}{\frac{1}{\sqrt{2n}}} \leq t\right) & =
 \frac{1}{\sqrt{2 \pi}} 
 \int_{ - t }^{\infty} e^{-\frac{x^2}{2}} dx
 \\
 & = \frac{1}{\sqrt{2 \pi}} 
 \int_{ - \infty }^{t} e^{-\frac{x^2}{2}} dx = \mathbb{P} \left(g \leq t\right),
\end{align*}
where $g \sim N(0,1)$.
\end{proof}


\section{Central Limit Theorems}
\label{section4}

In this section, we are going to prove the Central Limit Theorems \ref{Thm:CLTforlpballs} and \ref{Thm:CLTforCube}. We divide the section into two subsections. The first one is devoted to random vectors uniformly distributed on $l_p^n$-balls $(1 \leq p < \infty)$ or on their boundaries according to the cone measure, in order to prove Theorem \ref{Thm:CLTforlpballs}, and the second one is devoted to the case of the cube (Theorem \ref{Thm:CLTforCube}).

\subsection{CLT for $l_p^n-$balls}
Let $X^{(n)},Y^{(n)}$ be independent random vectors uniformly distributed on the boundary of $B_p^n$. By the Schechtman-Zinn representation (Lemma \ref{Lem:SchAndZinn}), 
\begin{align*}
X^{(n)} \overset{d}{=} \frac{G}{\Vert G \Vert _p} \; \text{ and } \; Y^{(n)} \overset{d}{=} \frac{G'}{\Vert G' \Vert _p},
\end{align*}
where $G=(g_1,...,g_n)$ and $G'=(g'_1,...,g'_n)$ are independent $p$-generalized Gaussian vectors on $\R^n$, that is, their components are $p$-generalized Gaussian random variables on $\R$.
Then, we can see the distance between random vectors uniformly distributed on the boundary of $B_p^n$, normalized by $n^{1/p - 1/2}$, as 
\begin{align*}
\dist( X^{(n)} , Y^{(n)}  ) & = n^{1/p - 1/2} \left\Vert \frac{G}{\Vert G \Vert_p} - \frac{G'}{\Vert G' \Vert_p} \right\Vert _2 
 = n^{1/p - 1/2} \left\Vert \frac{(g_1,...,g_n)}{\left(\sum _{i=1}^n \left| g_i  \right|^p \right)^{1/p}} - 
 \frac{(g'_1,...,g'_n)}{\left(\sum _{i=1}^n \left| g_i'  \right|^p \right)^{1/p}} \right\Vert _2
 \\
 & = \frac{n^{1/p - 1/2}}{\left(\sum _{i=1}^n \left| g_i  \right|^p \right)^{1/p}} \left\Vert (g_1,...,g_n) - 
 \frac{\left(\sum _{i=1}^n \left| g_i  \right|^p \right)^{1/p} (g'_1,...,g'_n)}{\left(\sum _{i=1}^n \left| g_i'  \right|^p \right)^{1/p}} \right\Vert _2
 \\
 & =  \frac{n^{1/p - 1/2}}{\left(\sum _{i=1}^n \left| g_i  \right|^p \right)^{1/p}} \left\Vert   \left( g_i - 
 \frac{\left(\sum _{i=1}^n \left| g_i  \right|^p  \right)^{1/p} }{\left(\sum _{i=1}^n \left| g_i'  \right|^p \right)^{1/p}} g_i' \right)_{i=1}^n \right\Vert _2.
\end{align*}
Note that by the Strong law of large numbers (SLLN), 
\begin{align*}
\frac{\left(\sum _{i=1}^n \left| g_i  \right|^p  \right)^{1/p} }{\left(\sum _{i=1}^n \left| g_i'  \right|^p \right)^{1/p}} \overset{a.s.}{\longrightarrow} 1 .
\end{align*}
Thus, the idea of the proof will be to first prove that $\dist( X^{(n)} , Y^{(n)} )$ and 
\begin{align}
\label{ExpSimplifCLTboundarypballs}
\frac{n^{1/p-1/2}}{\left(\sum _{i=1}^n \left| g_i  \right|^p \right)^{1/p}} \left\Vert   \left( g_i - 
   g_i' \right)_{i=1}^n \right\Vert _2,
\end{align}
have the same asymptotic distribution, and then prove a CLT for \eqref{ExpSimplifCLTboundarypballs}. This will imply that $\dist(X^{(n)} , Y^{(n)})$ satisfies the same CLT.


\begin{lemma}
\label{LemEquivalentCLTBoundarypball}
Let $1\leq p <\infty$ and let $g_i, g_i'$ with $i=1,...,n$, be independent $p-$generalized Gaussian random variables. Then,
\begin{align*}
\frac{n^{1/p - 1/2}}{\left(\sum _{i=1}^n \left| g_i  \right|^p \right)^{1/p}} \left\Vert   \left( g_i - 
 \frac{\left(\sum _{i=1}^n \left| g_i  \right|^p  \right)^{1/p} }{\left(\sum _{i=1}^n \left| g_i'  \right|^p \right)^{1/p}} g_i' \right)_{i=1}^n \right\Vert _2 \overset{d}{\sim} \frac{n^{1/p - 1/2}}{\left(\sum _{i=1}^n \left| g_i  \right|^p \right)^{1/p}} \left\Vert   \left( g_i - 
   g_i' \right)_{i=1}^n \right\Vert _2 .
\end{align*}
\end{lemma}

\begin{proof}
Let $G=(g_1,...,g_n)$ and $G'=(g'_1,...,g'_n)$ be independent $p$-generalized Gaussian vectors on $\R^n$.
Take $A_n := \frac{\Vert G \Vert _p}{\Vert G' \Vert _p} $ and $f(x) := \Vert G - xG' \Vert_2$, $\forall x \in \R$. By the mean value theorem, for almost every value of $G$ and $G'$ there exists some $\xi \in \R$ such that
\begin{align*}
f(1) - f(A_n) =  f'(\xi) (1- A_n).
\end{align*}
Note that $f'$ is differentiable for all $x$ except possibly when $G = xG'$. This event has probability $0$. Therefore, with probability $1$, $f$ is differentiable everywhere, and the mean value theorem applies.
Then,
\begin{align*}
\Vert G - G' \Vert_2 - \Vert G - A_n G' \Vert_2  = \left\langle \frac{G-\xi G'}{\Vert G - \xi G' \Vert} ,  (1- A_n) G' \right\rangle,
\end{align*}
for some $\xi \in \R$. Then, using Cauchy-Schwarz inequality, we obtain that
\begin{align*}
\left| \Vert G - G' \Vert_2 - \Vert G - A_n G' \Vert_2 \right|  \leq  | 1- A_n| \Vert G' \Vert_2.
\end{align*}
Equivalently,
\begin{align*}
\left| \; n^{-1/2}\left\Vert  G - A_n G' \right\Vert_2 - n^{-1/2} \left\Vert  G - G'  \right\Vert_2 \right| \leq n^{-1/2} | 1- A_n| \Vert G' \Vert_2.
\end{align*}
Let us bound this last expression. For the first factor, by the SLLN,
\begin{align*}
A_n = \frac{\Vert G \Vert _p}{\Vert G' \Vert _p}  \overset{a.s.}{\to} 1,
\end{align*}
therefore,
\begin{align*}
|1-A_n| \overset{a.s.}{\to} 0.
\end{align*}
For the second factor, again by the SLLN, there exists come constant $C:=\sqrt{M_p(2)}>0$ such that
\begin{align*}
n^{-1/2} \Vert G' \Vert _2  \overset{a.s.}{\longrightarrow} C .
\end{align*}
Then, it is clear that
\begin{align*}
\left| \; n^{-1/2}\left\Vert  G - A_n G' \right\Vert_2 - n^{-1/2} \left\Vert  G - G'  \right\Vert_2 \right| \overset{a.s.}{\to} 0.
\end{align*}
It follows that
\begin{align*}
 n^{-1/2}\left\Vert  G - A_n G' \right\Vert_2  \overset{d}{\sim} n^{-1/2} \left\Vert  G - G'  \right\Vert_2 .
\end{align*}
Therefore, since $n^{1/p} \Vert G \Vert _p^{-1} \overset{a.s.}{\to} M_p(p) ^{-1/p} >0$, by Slutsky Theorem we obtain that
\begin{align*}
 n^{1/p-1/2} \frac{1}{\Vert G \Vert_p} \left\Vert  G - A_n G' \right\Vert_2  \overset{d}{\sim} n^{1/p-1/2} \frac{1}{\Vert G \Vert_p} \left\Vert  G - G'  \right\Vert_2 .
\end{align*}
\end{proof}

Let us prove Thorem \ref{Thm:CLTforlpballs} in the case that $X^{(n)}, Y^{(n)} \sim \text{Unif} ( \partial B_p^ n)$.

\begin{proof}[Proof of Theorem \ref{Thm:CLTforlpballs} (boundary case)]
Let us define the function $f: \R^2 \to \R$ given by
\begin{align*}
f(x,y) =  x^{1/2} \frac{1}{y^{1/p}}.
\end{align*} 
First, note that
\begin{align*}
n^{1/p - 1/2} \frac{1}{\Vert G \Vert _p} \Vert G - G' \Vert _2 = f \left( \frac{1}{n} \sum _{i=1}^n \left( \left| g_i - g_i' \right|^2 , |g_i|^p \right)  \right).
\end{align*}
Take $S_{1,n} = \frac{1}{n} \sum _{i=1}^n  \left| g_i - g_i' \right|^2$ and $S_{2,n} = \frac{1}{n} \sum _{i=1}^n  \left| g_i  \right|^p$. Thus, by Lemma \ref{LemEquivalentCLTBoundarypball} we have that 
\begin{align*}
n^{1/p - 1/2} \Vert X^{(n)} - Y^{(n)} \Vert _2 \overset{d}{\sim} f\left( (S_{1,n} , S_{2,n}) \right).
\end{align*}
By the classical CLT, $\left( \sqrt{n} ( S_{1,n} - \E [S_{1,n}]), \sqrt{n} ( S_{2,n} - \E [S_{2,n}]) \right)$ converges in distribution to $N(0, \Sigma )$, where $\Sigma$ is the covariance matrix of $(S_{1,n},S_{2,n})$. Thus, applying the Delta method (Lemma \ref{DeltaMethod}) to $X_n = (S_{1,n} , S_{2,n})$ and the function $f$, we have that 
\begin{align*}
\sqrt{n} \left( f(X_n) - f\left( \E [X_n] \right) \right) \overset{d}{\to} N \left(0, J_f \Sigma J_f^T \right),
\end{align*}
where $J_f$ is the Jacobian matrix of $f$ and $\E[X_n] = (2M_p(2) , M_p(p))$.

Let us compute the covariance matrix 
\[
   \Sigma = 
  \left[ {\begin{array}{cc}
   \var[|g_1 - g_1'|^2] & \cov[|g_1 - g_1'|^2 , |g_1|^p] \\
   \cov[|g_1 - g_1'|^2 , |g_1|^p] &  \var[|g_1|^p] \\
  \end{array} } \right].
\]
For the first term,
\begin{align*}
\var \left[ \left| g_i - g_i' \right|^2   \right] & = \var \left[ | g_1 |^2 + | g_1' |^2
- 2 g_1 g_1'   \right]
 \\
 & = \E \left[ \left( | g_1 |^2 + | g_1' |^2 - 2 g_1 g_1' \right) ^2 \right] -
\left( \E \left[ | g_1 |^2 + | g_1' |^2 - 2 g_1 g_1'   \right] \right)^2
 \\
 & = \E \left[ |g_1|^4 + |g_1'|^4 + 6 |g_1|^2 |g_1'|^2 -4|g_1|^2 g_1g_1' -4|g_1'|^2 g_1g_1'   \right]
 - \left( 2 \E \left[ |g_1|^2  \right] \right)^2
 \\
 & = 2 \E \left[ |g_1|^4 \right] + 6 E\left[ |g_1|^2 \right] \E \left[ |g_1'|^2 \right] - \left( 2 \E \left[ |g_1|^2  \right] \right)^2
 \\
 & = 2 \E \left[ |g_1|^4 \right] + 2 \left( \E \left[ |g_1|^2  \right] \right)^2
 \\
 & = 2 M_p(4) + 2 (M_p(2))^2.
\end{align*}
For the fourth term,
\begin{align*}
\var \left[ \left| g_i\right|^p   \right] = \E \left[  | g_1 |^{2p} \right] -
\left( \E \left[ | g_1 |^p   \right] \right)^2 & = \frac{\Gamma \left( \frac{2p + 1}{p} \right)}{\Gamma \left( \frac{1}{p} \right)} - \left( \frac{\Gamma \left( \frac{p+1}{p} \right)}{\Gamma \left( \frac{1}{p} \right)} \right)^2 
 \\
  & = \frac{1}{p} = M_p(p).
\end{align*}
For the remaining terms, note that
\begin{align*}
\cov[\left| g_i - g_i'\right|^2 , \left| g_i\right|^p ] & = \E [\left| g_i - g_i'\right|^2 \left| g_i\right|^p  ] - \E [\left| g_i - g_i'\right|^2  ] \E [ \left| g_i\right|^p  ].
\end{align*}
Let us compute both summands. For the second one,
\begin{align*}
\E [\left| g_i - g_i'\right|^2] = 2\E [\left| g_1 \right|^2] = 2 \frac{\Gamma \left( \frac{3}{p} \right)}{\Gamma \left( \frac{1}{p} \right)}.
\end{align*}
and
\begin{align*}
\E [\left| g_i \right|^p] = \frac{\Gamma \left( \frac{p+1}{p} \right)}{\Gamma \left( \frac{1}{p} \right)} = \frac{1}{p}.
\end{align*}
For the first one,
\begin{align*}
\E [ \left| g_i - g_i'\right|^2  \left| g_i \right|^p] & = \E[ |g_1|^{p+2} + |g_1|^p|g_1'|^2 - 2g_1 g_1' |g_1|^p ]
 \\
  & =  \E[ |g_1|^{p+2}] +  \E[ |g_1|^{p}]  \E[ |g_1'|^{2}]
 \\
  & = \frac{\Gamma \left( \frac{p+3}{p} \right)}{\Gamma \left( \frac{1}{p} \right)} + \frac{\Gamma \left( \frac{p+1}{p} \right)}{\Gamma \left( \frac{1}{p} \right)} \frac{\Gamma \left( \frac{3}{p} \right)}{\Gamma \left( \frac{1}{p} \right)} = \frac{4\Gamma \left( \frac{p+3}{p} \right)}{3\Gamma \left( \frac{1}{p} \right)} .
\end{align*}
Therefore,
\begin{align*}
\cov[\left| g_i - g_i'\right|^2 , \left| g_i\right|^p ] & = \frac{4\Gamma \left( \frac{p+3}{p} \right)}{3\Gamma \left( \frac{1}{p} \right)} - 2 \frac{\Gamma \left( \frac{3}{p} \right)}{\Gamma \left( \frac{1}{p} \right)} \frac{1}{p}
 \\
 & = \frac{2 \Gamma \left( \frac{p+3}{p} \right)}{3\Gamma \left( \frac{1}{p} \right)} = \frac{2}{3} M_{p} (p+2).
\end{align*}
Thus, we obtain the covariance matrix
\[
   \Sigma = 
  \left[ {\begin{array}{cc}
   2 M_p(4) + 2 (M_p(2))^2 & \frac{2}{3} M_{p} (p+2) \\
   \frac{2}{3} M_{p} (p+2)  &  M_p(p) \\
  \end{array} } \right].
\]
For the last step, let us compute the Jacobian of $f: \R^2 \to \R$ given by
\begin{align*}
f(x,y) =  x^{1/2} \frac{1}{y^{1/p}}  .
\end{align*} 
It is clear that 
\[
   J_f (x,y) = 
  \left[ {\begin{array}{cc}
   \frac{1}{2\sqrt{x}} \frac{1}{y^{1/p}}  & \frac{-\sqrt{2}}{p} \frac{\sqrt{x}}{y^{1+ 1/p}}  \\
  \end{array} } \right] .
\]
Thus, the Jacobian matrix of $f$ evaluated at $\E[X_n] = \left( 2 M_p(2) , 1/p \right)$ is
\[
   J_g = 
  \left[ {\begin{array}{cc}
   \frac{p^{1/p} }{2 \sqrt{2} \sqrt{M_p(2)}}  & - \sqrt{2}  p^{1/p} \sqrt{M_p(2)} \\
  \end{array} } \right] .
\]
Then, applying elementary computations, one can check that the vaue of $\sigma_p ^2$ is given by
\begin{align*}
\sigma_p ^2 = J_g \Sigma J_g ^T = \frac{p^{2/p} \left( M_p(4) + (M_p(2))^2 \right)}{4 M_p(2)}.
\end{align*}
\end{proof}

Now, let us work on the case where $X^{(n)},Y^{(n)}$ are random vectors uniformly distributed on $B_p^n$. By the Schechtman-Zinn representation (Lemma \ref{Lem:SchAndZinn}), 
\begin{align*}
X^{(n)} \overset{d}{=} U^{1/n} \frac{G}{\Vert G \Vert _p} \; \text{ and } \; Y^{(n)} \overset{d}{=} U'^{1/n} \frac{G'}{\Vert G' \Vert _p},
\end{align*}
where $G=(g_1,...,g_n)$ and $G'=(g'_1,...,g'_n)$ are independent $p$-generalized Gaussian vectors on $\R^n$, and $U,U'$ are uniformly distributed in the interval $[0,1]$ and independent of $G$ and $G'$.
Then, similarly to the previous case, we can see the distance between random vectors uniformly distributed on $B_p^n$, normalized by $n^{1/p - 1/2}$, as
\begin{align*}
n^{1/p - 1/2}\Vert X - Y \Vert _2 & = \left\Vert U^{1/n} \frac{G}{\Vert G \Vert_p} -U'^{1/n} \frac{G'}{\Vert G' \Vert_p} \right\Vert _2 
 \\
 & = n^{1/p - 1/2} \frac{U^{1/n}}{\Vert G \Vert _p} \left\Vert G-  \frac{U'^{1/n}}{U^{1/n}} \frac{\Vert G \Vert_p}{\Vert G' \Vert_p} G' \right\Vert _2 ,
\end{align*}
where there is an extra $U^{1/n}$ factor.
Note that since $U^{1/n}, U'^{1/n} \overset{a.s.}{\to} 1$ and $U^{1/n}, U'^{1/n}>0$ almost surely, we have that
\begin{align*}
\frac{U'^{1/n}}{U^{1/n}} \overset{a.s.}{\to} 1 .
\end{align*}
Thus, again we will first prove that $\dist( X^{(n)} , Y^{(n)} )$ and 
\begin{align}
\label{ExpSimplifCLTpballs}
\frac{n^{1/p-1/2}}{\left(\sum _{i=1}^n \left| g_i  \right|^p \right)^{1/p}} U^{1/n} \left\Vert   \left( g_i -  g_i' \right)_{i=1}^n \right\Vert _2,
\end{align}
have the same asymptotic distribution, and then prove a CLT for \eqref{ExpSimplifCLTpballs}. This will imply that $\dist(X^{(n)} , Y^{(n)})$ satisfies the same CLT.


\begin{lemma}
\label{LemEquivalentCLTpball}
Let $1\leq p <\infty$ and let $g_i, g_i'$ with $i=1,...,n$, be independent $p-$generalized Gaussian random variables and $U,U'$ be independent uniform random variables on the interval $[0,1]$. Then,
\begin{align*}
n^{1/p - 1/2} \frac{U^{1/n}}{\Vert G \Vert _p} \left\Vert G-  \frac{U'^{1/n}}{U^{1/n}} \frac{\Vert G \Vert_p}{\Vert G' \Vert_p} G' \right\Vert _2 \overset{d}{\sim} n^{1/p - 1/2} \frac{U^{1/n}}{\Vert G \Vert _p} \left\Vert G-  G' \right\Vert _2 .
\end{align*}
\end{lemma}

\begin{proof}
The argument is the same as in the proof of Lemma \ref{LemEquivalentCLTBoundarypball}, replacing
\[
A_n=\frac{\|G\|_p}{\|G'\|_p}
\quad\text{by}\quad
A_n=\frac{U'^{1/n}}{U^{1/n}}\frac{\|G\|_p}{\|G'\|_p}.
\]
In particular, defining $f(x)=\|G-xG'\|_2$, the same mean value theorem argument yields
\[
\big|\,n^{-1/2}\|G-A_nG'\|_2-n^{-1/2}\|G-G'\|_2\,\big|
\le |1-A_n|\,n^{-1/2}\|G'\|_2
\qquad\text{a.s.}
\]
By the SLLN, $\|G\|_p/\|G'\|_p\to 1$ a.s. and $n^{-1/2}\|G'\|_2\to \sqrt{M_p(2)}$ a.s.
Moreover, $U^{1/n}\to 1$ and $U'^{1/n}\to 1$ a.s., hence $A_n\to 1$ a.s., and therefore $|1-A_n|\to 0$ a.s..
Consequently, the right-hand side converges to $0$ a.s., which implies that
\[
n^{-1/2}\|G-A_nG'\|_2 \overset{d}{\sim} n^{-1/2}\|G-G'\|_2.
\]
Finally, since $n^{1/p}\|G\|_p^{-1}\to M_p(p)^{-1/p}>0$ a.s., an application of Slutsky's theorem gives
\[
n^{1/p-1/2}\frac{U^{1/n}}{\|G\|_p}\|G-A_nG'\|_2
\overset{d}{\sim}
n^{1/p-1/2}\frac{U^{1/n}}{\|G\|_p}\|G-G'\|_2.
\]
\end{proof}

Now, we are ready to prove Theorem \ref{Thm:CLTforlpballs}.
\begin{proof}[Proof of Theorem \ref{Thm:CLTforlpballs} ($B_p^n$-case)]
Let $X^{(n)}, Y^{(n)} \sim \text{Unif} (B_p^n)$ be independent. By Lemma \ref{LemEquivalentCLTpball}, it is enough to prove a CLT for \eqref{ExpSimplifCLTpballs}.
In particular, notice that 
\begin{align*}
U^{1/n} = e^{\frac{1}{n} \log (U)} \overset{d}{=} e^{-\frac{E}{n}},
\end{align*}
where $E$ is an exponential random variable with mean $1$. Therefore, using the series expansion of the exponential,
\begin{align*}
 e^{-\frac{E}{n}} = 1- \frac{E}{n} + R_n,
\end{align*}
where $R_n$ denotes the Taylor reminder term, and $|R_n| \leq E^2 / 2n^2$. Then,
\begin{align*}
e^{-\frac{E}{n}} = 1- \frac{E}{n} + O _E \left( \frac{1}{n^2} \right),
\end{align*}
where the Landau symbol $O_E$ depends on the realization of the random variable $E$. Since $E^2 < \infty$ almost surely, the bound above also implies that 
\begin{align*}
e^{-\frac{E}{n}} = 1- \frac{E}{n} + O \left( \frac{1}{n^2} \right)
\end{align*}
almost surely. Therefore, we have that 
\begin{align*}
\frac{U^{1/n}}{\Vert G \Vert _p} \left\Vert G -   G' \right\Vert _2 & \overset{d}{=} \frac{\left( 1- \frac{E}{n} + O \left( \frac{1}{n^2} \right)  \right)}{\Vert G \Vert _p} \left\Vert G -   G' \right\Vert _2
 \\
 & \overset{d}{\longrightarrow} \frac{1}{\Vert G \Vert _p} \left\Vert G -   G' \right\Vert _2 .
\end{align*}

This last expression coincides with the case where $X^{(n)}, Y^{(n)}$ are uniformly distributed on the boundary of $B_p^n$. Thus, the Theorem follows.
\end{proof}


\subsection{CLT for the cube}
\label{Subsection:CLTCube}

Let $X^{(n)},Y^{(n)}$ be independent random vectors uniformly distributed on $B_\infty ^n$. Then $X^{(n)} \overset{d}{=} (u_1,...,u_n)$ and $Y^{(n)} \overset{d}{=} (u'_1,...,u'_n)$ where $u_i, u'_i$ are independent random variables uniformly distributed on the interval $[-1,1]$. Thus, the proof of the CLT for this case is a straightforward computation.



Let us first prove that the case where $X^{(n)},Y^{(n)}$ are uniformly distributed on the cube, is asymptotically distributed as the case where $X^{(n)},Y^{(n)}$ are uniformly distributed on the boundary of the cube.

\begin{proof}[proof of Theorem \ref{Thm:CLTforCube}]
Let $X^{(n)},Y^{(n)}$ be random vectors uniformly distributed on $B_\infty ^n$. Then $X^{(n)} \overset{d}{=} (u_1,...,u_n)$ and $Y^{(n)} \overset{d}{=} (u'_1,...,u'_n)$ where $u_i, u'_i$ are independent random variables uniformly distributed on the interval $[-1,1]$.

Note that
\begin{align*}
n^{-1/2}\Vert X^{(n)} - Y^{(n)} \Vert _2 & = \left\Vert  (u_1,...,u_n) - (u_1',...,u_n') \right\Vert
 \\
 & = \left( \frac{1}{n} \sum _{i=1}^n \left| u_i - u_i' \right|^2 \right)^{1/2} .
\end{align*}
For each $n\in \N$, we take
\begin{align}
\label{SnCLTCube}
S_n := \frac{1}{n} \sum _{i=1}^n \left| u_i - u_i' \right|^2.
\end{align}
By the classical CLT, it is clear that $\sqrt{n}\left( S_n - \mu \right) \to N (0 , \sigma ^2)$, where 
\begin{align*}
\mu := \E[ |u_1 - u_1' |^2 ] = \frac{2}{3}
\end{align*}
and
\begin{align*}
\sigma^2 := \var \left[ |u_1-u_1'|^2 \right] = \E\left[ |u_1-u_1'|^4 \right] - \left( \E\left[ |u_1-u_1'|^2 \right]\right)^2 = 
\frac{2^4}{15} - \left(\frac{2}{3}\right)^2 = \frac{28}{45}.
\end{align*}
Thus, taking the function $f(x) = \sqrt{x}$, we have that
\begin{align*}
n^{-1/2}\Vert X^{(n)} - Y^{(n)} \Vert _2 & = f\left( S_n  \right).
\end{align*}
Applying the Delta method (\ref{DeltaMethod}), we have that
\begin{align*}
\sqrt{n} \left( n^{-1/2}\Vert X^{(n)} - Y^{(n)} \Vert _2 - f (\mu ) \right)  \overset{d}{\to} N \left( 0 , \sigma ^2 (f'(\mu))^2  \right) = N\left( 0 , \frac{7}{30} \right).
\end{align*}


Now, let $ \tilde{X}^{(n)}, \tilde{Y}^{(n)}$ be random vectors uniformly distributed on $\partial B_\infty ^n$. Note that any point on the boundary of the hypercube has one coordinate either $1$ or $-1$, and the rest of coordinates are uniform random variables on the interval $[-1,1]$. Therefore, we take $X^{(n)} = (u_1, ... , u_n)$ and $Y^{(n)} = (u'_1, ... , u'_n)$ as before, and we define $\tilde{X}^{(n)} = (\tilde{u}_1, ... , \tilde{u}_n)$ as
\[ 
\tilde{u}_i = 
 \begin{cases}
  u_i, & \text{if }i \neq K_X \\
  S_X & \text{if }i=K_X \\
 \end{cases}
\]
where $K_X \sim \text{Unif}(\lbrace 1 , ... , n \rbrace)$ and $S_X \sim \text{Unif}(\lbrace -1 , 1 \rbrace)$, and $\tilde{Y}^{(n)} = (\tilde{u'}_1, ... , \tilde{u'}_n)$ as
\[ 
\tilde{u'}_i = 
 \begin{cases}
  u'_i, & \text{if }i \neq K_Y \\
  S_Y & \text{if }i=K_Y \\
 \end{cases}
\]
where $K_Y \sim \text{Unif}(\lbrace 1 , ... , n \rbrace)$ and $S_Y \sim \text{Unif}(\lbrace -1 , 1 \rbrace)$.
For each $n\in \N$, we take
\begin{align*}
S_n := \frac{1}{n} \sum _{i=1} ^n |u_i - u'_i |^2, \text{ and }  \tilde{S}_n := \frac{1}{n} \sum _{i=1} ^n |\tilde{u}_i - \tilde{u'}_i |^2 .
\end{align*}
Note that $S_n$ differs from $\tilde{S_n}$ in at most two summands, for which
\begin{align*}
\left| (u_i - u'_i )^2 - (\tilde{u}_i - \tilde{u'}_i )^2 \right| \leq 8.
\end{align*}
since $|u_i|, |u'_i|, |\tilde{u}_i|, |\tilde{u'}_i| \leq 1$, for any $i\in \lbrace 1, ... , n\rbrace$. Then, it follows that
\begin{align}
\label{Ineq:SnandSnTilde}
|S_n - \tilde{S}_n | = \left| \frac{1}{n} \sum _{i=1}^n \left( (u_i - u'_i )^2 - (\tilde{u}_i - \tilde{u'}_i )^2  \right) \right| \leq 2 \frac{8}{n} = \frac{16}{n}
\end{align}
almost surely. Note that, 
\begin{align*}
\left| n^{-1/2} \Vert X^{(n)} - Y^{(n)} \Vert _2 - n^{-1/2} \Vert \tilde{X}^{(n)} - \tilde{Y}^{(n)} \Vert _2  \right| & = \left| \sqrt{S_n} - \sqrt{\tilde{S}_n} \right|
 \\
  & = \frac{\left| S_n - \tilde{S}_n \right|}{ \sqrt{S_n} + \sqrt{\tilde{S}_n} }
\end{align*}
almost surely, since $S_n, \tilde{S}_n >0 $ with probability $1$. 
Therefore, since $S_n \overset{a.s.}{\to} \E [S_n] >0$ and $\tilde{S_n} \overset{a.s.}{\to} \E [S_n]$, there exists some $c>0$ such that $ \sqrt{S_n} + \sqrt{\tilde{S}_n} > c$ almost surely, and using \eqref{Ineq:SnandSnTilde}, we obtain that
\begin{align*}
\left| n^{-1/2} \Vert X^{(n)} - Y^{(n)} \Vert _2 - n^{-1/2} \Vert \tilde{X}^{(n)} - \tilde{Y}^{(n)} \Vert _2  \right| 
 \leq \frac{16/n}{c} \overset{a.s.}{\to} 0
\end{align*}
which implies that
\begin{align*}
n^{-1/2} \Vert X^{(n)} - Y^{(n)} \Vert _2 \overset{d}{\sim} n^{-1/2} \Vert \tilde{X}^{(n)} - \tilde{Y}^{(n)} \Vert _2.
\end{align*}
Therefore, $n^{-1/2} \Vert \tilde{X}^{(n)} - \tilde{Y}^{(n)} \Vert _2$ satisfies the same CLT as $n^{-1/2} \Vert X^{(n)} - Y^{(n)} \Vert _2$.
\end{proof}










\section{Large Deviation Principles}
\label{section5}

In this section, we are going to prove the Large Deviation Principles results. We divide the section into two subsections. The first one is devoted to random vectors uniformly distributed on $l_p^n$-balls $(2 \leq p < \infty)$ in order to prove Theorems \ref{ThmLDPBoundarypgeq2} and \ref{ThmLDPpgeq2}, and the second one is devoted to the case of the cube (Theorem \ref{ThmLDPBoundaryCube}).



\subsection{LDP for $l_p^n-$balls, $p\geq 2$}


As in the previous section, instead of proving that an LDP is satisfied by $n^{1/p - 1/2} \Vert X^{(n)} - Y^{(n)} \Vert _2$, we show that it is satisfied by
\begin{align*}
W_n : = n^{1/p - 1/2} \frac{1}{\left(\sum _{i=1}^n \left| g_i  \right|^p \right)^{1/p}} \left(  \sum_{i=1}^n \left| g_i - g_i' \right|^2 \right)^{1/2},
\end{align*} 
where $(g_i)_{i=1} ^n$ and $(g'_i)_{i=1} ^n$ are independent $p$-generalized Gaussian random variables.
For that, we need to first prove that these two expressions are exponentially equivalent. In that case, by Theorem \ref{LemExponentialEquivalence}, both will satisfy the same LDP with speed $n$ and the same rate function.


Let us first see that the two expressions are exponentially equivalent.
\begin{lemma}
\label{LemExpEquivalentpgeq2Boundary}
Let $p\geq 2$, and let $G=(g_1,...,g_n), G'=(g_1',...,g_n')$ be independent generalized $p-$gaussian vectors, and let $U,U'$ be independent uniform random variables on the interval $[0,1]$, independent of $G$ and $G'$. Let
\begin{align*}
A_n = \frac{\Vert G \Vert _p}{\Vert G' \Vert _p} \; \text{ or } \; A_n = \frac{U^{1/n}}{U'^{1/n}} \frac{\Vert G \Vert _p}{\Vert G' \Vert _p}.
\end{align*} 
Then,
\begin{align*}
\frac{n^{1/p - 1/2}}{\Vert G \Vert _p} \left\Vert  G - A_n G' \right\Vert_2
\end{align*}
is exponentially equivalent to
\begin{align*}
\frac{n^{1/p - 1/2}}{\Vert G \Vert _p} \left\Vert  G - G' \right\Vert_2.
\end{align*}
Equivalently, for every $\delta >0$, 
\begin{align*}
\limsup _{n \to \infty}  \frac{1}{n} \log \mathbb{P} \left[ \left|  \frac{n^{1/p - 1/2}}{\Vert G \Vert _p} \left\Vert  G - A_n G' \right\Vert_2 - 
\frac{n^{1/p - 1/2}}{\Vert G \Vert _p} \left\Vert  G - G' \right\Vert_2
 \right|  > \delta \right] = -\infty .
\end{align*}
\end{lemma}

\begin{proof}

Let $f(x)=\|G-xG'\|_2$. Arguing as in the proof of Lemma \ref{LemEquivalentCLTBoundarypball}, we obtain that
\begin{align*}
\left| \Vert G - G' \Vert_2 - \Vert G - A_n G' \Vert_2 \right|  \leq  | 1- A_n| \Vert G' \Vert_2
\end{align*}
almost surely.
Now, note that for any $\delta >0$ and $\epsilon \in (0,1)$, we have that
\begin{align*}
& \mathbb{P} \left[ \left|  \frac{n^{1/p - 1/2}}{\Vert G \Vert _p} \left\Vert  G - A_n G' \right\Vert_2 - 
\frac{n^{1/p - 1/2}}{\Vert G \Vert _p} \left\Vert  G - G' \right\Vert_2
 \right|  > \delta \right] 
 \\
 = &  \mathbb{P} \left[  \frac{n^{1/p - 1/2}}{\Vert G \Vert _p} 
 \left|  \left\Vert  G - A_n G' \right\Vert_2 - 
 \left\Vert  G - G' \right\Vert_2
 \right|  > \delta \right] 
 \\
\leq &  \mathbb{P} \left[ \left|  \frac{n^{1/p}}{\Vert G \Vert _p}
 | 1- A_n | \frac{  \left\Vert G'  \right\Vert _2   }{n^{1/2}} 
 \right|  > \delta \right]
 \\
\leq  &  \mathbb{P} \left[   \frac{n^{1/p}}{\Vert G \Vert _p}   > \alpha   \right]
 +  \mathbb{P} \left[ \left|  1-A_n \right| > \beta   \right]
 + \mathbb{P} \left[   \frac{  \left\Vert G'  \right\Vert_2 }{n^{1/2}}   \geq  \gamma   \right]
 \\
= &  \mathbb{P} \left[   \frac{\Vert G \Vert _p ^p}{n}   < \frac{1}{\alpha^p}   \right]
 +  \mathbb{P} \left[ \left|  1-A_n \right| > \beta   \right]
 + \mathbb{P} \left[   \frac{  \left\Vert G'  \right\Vert_2 ^2   }{n}   \geq \gamma^2   \right] .
\end{align*}
where $\alpha = 1/ M_p(p)^{1/p} + \epsilon$, $\gamma = (M_p(2) + \epsilon)^{1/2}$ and $\beta = \delta / ( \alpha \gamma ) $. In particular, noticing how $\alpha \gamma $ increases with $\epsilon$, we can take $\epsilon>0$ big enough such that $\beta \in (0, 1/2]$.


\quad

\quad


Let us bound these three terms. For the first term, by Cramer's theorem,
\begin{align*}
\mathbb{P} \left[ \left|  \frac{\Vert G \Vert _p ^p}{n} \right|  < \frac{1}{\alpha ^p}   \right]  & = \mathbb{P} \left[   \frac{1}{n} \sum _{i=1}^n \left| g_i  \right|^p   <  \frac{M_p(p)}{\left( 1 + \epsilon M_p(p)^{1/p} \right)^p}   \right] \leq e^{-c_1 n}
\end{align*}
for some $c_1>0$.


For the third term, by Cramer's theorem, 
\begin{align*}
\mathbb{P} \left[ \left|  \frac{  \left\Vert G'  \right\Vert_2 ^2   }{n} \right|  \geq  \gamma  ^2 \right] = \mathbb{P} \left[ \left|  \frac{  \left\Vert G'  \right\Vert_2 ^2   }{n} \right|  \geq  M_p(2) + \epsilon   \right] \leq e^{-c_2 n}
\end{align*}
for some $c_2 >0$.




For the second term, we differenciate two cases of $A_n$. If $A_n = \frac{\Vert G \Vert _p}{\Vert G' \Vert _p}$, we have that for every $\delta >0$, and taking a suitable $\epsilon>0$ such that $\beta \in (0, 1/2]$,
\begin{align*}
& \mathbb{P} \left( \left| \frac{\left(\sum _{i=1}^n \left| g_i  \right|^p  \right)^{1/p} }{\left(\sum _{i=1}^n \left| g_i'  \right|^p \right)^{1/p}} - 1 \right| > \beta \right)  
 \\
 & \leq \mathbb{P} \left(  \frac{ \left( \frac{1}{n} \sum _{i=1}^n \left| g_i  \right|^p  \right)^{1/p} }{\left(\frac{1}{n}\sum _{i=1}^n \left| g_i'  \right|^p \right)^{1/p}} < 1- \beta \right) + 
 \mathbb{P} \left(  \frac{ \left( \frac{1}{n} \sum _{i=1}^n \left| g_i  \right|^p  \right)^{1/p} }{\left(\frac{1}{n}\sum _{i=1}^n \left| g_i'  \right|^p \right)^{1/p}} > 1+ \beta \right)
  \\
  & \leq \mathbb{P} \left( \frac{1}{n} \sum _{i=1}^n \left| g_i  \right|^p  < (1- \beta)^{p/2} M_p(p) \right)
  + \mathbb{P} \left( \frac{1}{n} \sum _{i=1}^n \left| g_i'  \right|^p  > (1- \beta)^{-p/2} M_p(p) \right)
  \\
  & + \mathbb{P} \left( \frac{1}{n} \sum _{i=1}^n \left| g_i  \right|^p  > (1+ \beta)^{p/2} M_p(p) \right)
  + \mathbb{P} \left( \frac{1}{n} \sum _{i=1}^n \left| g_i'  \right|^p  < (1+ \beta)^{-p/2} M_p(p) \right).
\end{align*} 
By Cramer's theorem, the four terms decay exponentially with speed $n$.

In the second case, if $A_n = \frac{U^{1/n}}{U'^{1/n}} \frac{\Vert G \Vert _p}{\Vert G' \Vert _p}$, note that for every $\delta >0$, and taking a suitable $\epsilon>0$ such that $\beta \in (0, 1/2]$,
\begin{align*}
& \mathbb{P} \left( \left| \frac{U^{1/n}}{U'^{1/n}} \frac{\left(\sum _{i=1}^n \left| g_i  \right|^p  \right)^{1/p} }{\left(\sum _{i=1}^n \left| g_i'  \right|^p \right)^{1/p}} - 1 \right| > \beta \right)  
 \\
 & \leq \mathbb{P} \left( \frac{U^{1/n}}{U'^{1/n}} \frac{ \left( \frac{1}{n} \sum _{i=1}^n \left| g_i  \right|^p  \right)^{1/p} }{\left(\frac{1}{n}\sum _{i=1}^n \left| g_i'  \right|^p \right)^{1/p}} < 1- \beta \right) + 
 \mathbb{P} \left(  \frac{U^{1/n}}{U'^{1/n}} \frac{ \left( \frac{1}{n} \sum _{i=1}^n \left| g_i  \right|^p  \right)^{1/p} }{\left(\frac{1}{n}\sum _{i=1}^n \left| g_i'  \right|^p \right)^{1/p}} > 1+ \beta \right)
  \\
  & \leq \mathbb{P} \left( \frac{1}{n} \sum _{i=1}^n \left| g_i  \right|^p  < (1- \beta)^{p/4} M_p(p) \right)
  + \mathbb{P} \left( \frac{1}{n} \sum _{i=1}^n \left| g_i'  \right|^p  > (1- \beta)^{-p/4} M_p(p) \right)
  \\
  & + \mathbb{P} \left( \frac{1}{n} \sum _{i=1}^n \left| g_i  \right|^p  > (1+ \beta)^{p/4} M_p(p) \right)
  + \mathbb{P} \left( \frac{1}{n} \sum _{i=1}^n \left| g_i'  \right|^p  < (1+ \beta)^{-p/4} M_p(p) \right)
 \\
  & + \mathbb{P} \left( U^{1/n}  < (1 - \beta)^{p/4} \right) + 
   \mathbb{P} \left( U'^{1/n}  > (1 - \beta)^{-p/4} \right)
  \\
   & + \mathbb{P} \left( U^{1/n}  > (1 + \beta)^{p/4} \right) + 
   \mathbb{P} \left( U'^{1/n}  < (1 + \beta)^{-p/4} \right).
\end{align*} 
Applying Cramer's theorem as before, the first four terms decay exponentially with speed $n$. Since
\begin{align*}
\mathbb{P} \left( U^{1/n} < x \right) = e^{-n |\log (x) |}, \text{ if }x \in (0,1)
\end{align*}
and
\begin{align*}
\mathbb{P} \left( U^{1/n} > x \right) = 0, \text{ if }x \geq 1,
\end{align*}
the remaining four terms either decay exponentially with speed $n$ or their value is $0$.
\end{proof}



Now, we are ready to prove both LDP's.
\begin{proof}[Proof of Theorem \ref{ThmLDPBoundarypgeq2}]
As before, for $n \in \N$, we recall the random vector 
\begin{align*}
S_n = \left( S_{1,n} , S_{2,n} \right) =
 \left( \frac{1}{n} \sum _{i=1}^n  \left| g_i - g_i' \right|^2 , 
 \frac{1}{n} \sum _{i=1}^n  \left| g_i  \right|^p \right).
\end{align*}
Then, for $t = (t_1,t_2) \in \R^2$,
\begin{align*}
\Lambda (t_1,t_2) & = \log \E \left[ e^{\langle t,( |g_1 - g_1'|^2 ,|g_1|^p ) \rangle} \right]
 \\
  & = \log \int _{-\infty} ^{\infty} \int _{-\infty} ^{\infty} e^{t_1 |x-y|^2 + t_2|y|^p    
  - \frac{|x|^p}{p} - \frac{|y|^p}{p} } \frac{1}{4p^{2/p} \left( \Gamma(1 + 1/p) \right)^2} dx dy .
\end{align*}

Note that if $t_1 \geq 0$, using the inequality $|x-y|^2 \leq 2(|x|^2 + |y|^2)$, we obtain that
\begin{align*}
\Lambda (t_1,t_2) & = \log \int _{-\infty} ^{\infty} \int _{-\infty} ^{\infty} e^{t_1 |x-y|^2 + t_2|y|^p    
  - \frac{|x|^p}{p} - \frac{|y|^p}{p} } \frac{1}{4p^{2/p} \left( \Gamma(1 + 1/p) \right)^2} dx dy
  \\
  & \leq \log \int _{-\infty} ^{\infty} \int _{-\infty} ^{\infty} e^{
  2t_1 |x|^2 + 2t_1|y|^2 + t_2|y|^p    
  - \frac{|x|^p}{p} - \frac{|y|^p}{p} } \frac{1}{4p^{2/p} \left( \Gamma(1 + 1/p) \right)^2} dx dy
  \\
  & = \log \int _{-\infty} ^{\infty}  e^{
  2t_1 |x|^2 - \frac{|x|^p}{p}    } dx
   \int _{-\infty} ^{\infty} e^{2t_1|y|^2 + \left(t_2 - \frac{1}{p} \right)|y|^p   } 
    \frac{1}{4p^{2/p} \left( \Gamma(1 + 1/p) \right)^2} dy,
\end{align*}
and if $t_1 \leq 0$, then
\begin{align*}
\Lambda (t_1,t_2) & = \log \int _{-\infty} ^{\infty} \int _{-\infty} ^{\infty} e^{t_1 |x-y|^2 + t_2|y|^p    
  - \frac{|x|^p}{p} - \frac{|y|^p}{p} } \frac{1}{4p^{2/p} \left( \Gamma(1 + 1/p) \right)^2} dx dy
  \\
  & \leq \log \int _{-\infty} ^{\infty} \int _{-\infty} ^{\infty} e^{
    t_2|y|^p    
  - \frac{|x|^p}{p} - \frac{|y|^p}{p} } \frac{1}{4p^{2/p} \left( \Gamma(1 + 1/p) \right)^2} dx dy
  \\
  & = \log \int _{-\infty} ^{\infty}  e^{
   - \frac{|x|^p}{p}    } dx
   \int _{-\infty} ^{\infty} e^{  \left(t_2 - \frac{1}{p} \right)|y|^p   } 
    \frac{1}{4p^{2/p} \left( \Gamma(1 + 1/p) \right)^2} dy.
\end{align*}
Notice that if $p > 2$, its domain contains the set $(\R \times (-\infty,1/p)$, where the origin is an interior point. And if $p = 2$, then the domain contains the set $\lbrace (t_1 , t_2) \in \mathbb{R}^2 \; : \; 2t_1 < 1/p , 2t_1 + t_2 < 1/p , t_2 < 1/p \rbrace$, where again the origin is an interior point. Therefore, by Cramer's theorem the sequence $(S_n)_{n\in\N}$ satisfies an LDP in $\R^2$ with speed $n$ and good rate function $\Lambda ^*$.

Next, we define the continuous function $f(x,y) = x^{1/2} y^{-1/p}$.
As stated before, for each $n\in \N$, $W_n \overset{d}{=} f(S_n)$. Thus, by the contraction principle, the random sequence $W = (W_n)_{n \in \N}$ satisfies an LDP on $\R$ with speed $n$ and good rate function 

\[   
I_W(z) = 
     \begin{cases}
       \displaystyle{\inf_{ \substack{x \geq 0, y> 0 \\ x^{1/2}y^{-1/p} = z}} \Lambda ^* (x,y)}  &\quad\text{: }z \geq 0\\
       + \infty &\quad\text{: }z < 0 . \\ 
     \end{cases}
\]
\end{proof}

\begin{proof}[Proof of Theorem \ref{ThmLDPpgeq2}]
In the case of $X^{(n)}, Y^{(n)} \sim \text{Unif}(B_p^n)$, then by Lemma \ref{LemExpEquivalentpgeq2Boundary}, $n^{1/p - 1/2} \Vert X^{(n)} - Y^{(n)} \Vert _2$ satisfy an LDP with the same speed and same rate function as
\begin{align*}
 n^{1/p - 1/2} \frac{U^{1/n}}{\left(\sum _{i=1}^n \left| g_i  \right|^p \right)^{1/p}} \left(  \sum_{i=1}^n \left| g_i - g_i' \right|^2 \right)^{1/2} .
\end{align*} 
Then, let us define the random variables $V_n:=(U^{1/n} , W_n)$, where $W_n$ is defined as in Theorem \ref{ThmLDPBoundarypgeq2}. As proved in Lemma 3.3 of \cite{GKR}, $U:=(U^{1/n})_{n \in \N}$ satisfy an LDP on $\R$ with speed $n$ and rate function 
\[   
I_U(z) = 
     \begin{cases}
       - \log (z) &\quad\text{: }z \in (0, 1]\\
       + \infty &\quad\text{: otherwise}. \\ 
     \end{cases}
\]
As $U^{1/n}$ and $W_n$ are independent, the sequence $V:=(V_n)_{n \in \N}$ satisfies an LDP on $\R^2$ with speed $n$ and good rate function
\begin{align*}
I_V(z_1,z_2) := I_U(z_1) + I_W(z_2), \quad (z_1,z_2)\in \R^2.
\end{align*}
Finally, applying once again the contraction principle, now to the function $f(x,y) = xy$, we conclude that the sequence of random variables 
\begin{align*}
n^{1/p - 1/2}\Vert X^{(n)}-Y^{(n)}\Vert _2 \overset{d}{=} U^{1/n} W_n 
\end{align*}
satisfies and LDP on $\R$ with speed $n$ and good rate function 
\[   
I_{\Vert X-Y \Vert _2}(z) = \inf _ {z=z_1 z_2} I_V(z_1,z_2) = 
     \begin{cases}
       \displaystyle{\inf_{ \substack{z_1 \geq 0, z_2 \geq 0 \\ z=z_1z_2}} I_V (z_1,z_2)}  &\quad\text{: }z \geq 0\\
       + \infty &\quad\text{: otherwise}. \\ 
     \end{cases}
\]
\end{proof}


\subsection{LDP for $\infty$-balls }

In this section, we provide a proof of the LDP for the cube.

\begin{proof}[Proof of Theorem \ref{ThmLDPBoundaryCube}]
Using the same notation as in Section \ref{Subsection:CLTCube}, let $X_i = |u_i-u_i'|^2$ and $S_n = \frac{1}{n} \sum_{i=1}^n X_i$. By Cramer's theorem, we have that the sequence $(S_n)_{n\in\N}$ satisfies an LDP on $\R$ with speed $n$ and good rate function $\Lambda^*$, where
\begin{align*}
\Lambda (t) = \log E \left[ e^{X_i t} \right] = \log \left[ 2 \int _{0} ^2 e^{t x^2 } \left( 1- \frac{x}{2} \right) dx \right].
\end{align*}
Now, using the contraction principle with the function $f(x) = \sqrt{x}$, one can see that $f(S_n) \overset{d}{=} \text{dist}(X^{(n)} , Y^{(n)}) $ satisfies a LDP with the same speed and good rate function $\Lambda^* \circ f^{-1}$. That is,
\[   
I_{\text{dist}(X^{(n)} , Y^{(n)})}(z)  = 
     \begin{cases}
       \displaystyle{\inf_{ \substack{x \geq 0 \\ z = \sqrt{x}}} \Lambda^* (x)}  &\quad\text{: }z \geq 0\\
       + \infty &\quad\text{: otherwise}. \\ 
     \end{cases}
\]
\end{proof}




\begin{thebibliography}{99} % don't worry about the 99

\bibitem{VB}{\sc I. B\'ar\'any and V. Vu} \textit{Central limit theorems for Gaussian polytopes.} Ann. Probab., 35(4):1593-1621, 2007.

\bibitem{DF}{\sc P. Diaconis and D. Freedman.} \textit{A dozen de Finetti-style results in search of a theory.} Annales de l'I.H.P. Probabilit\'es et statistiques, Volume 23 (1987) no. S2, pp. 397-423.

\bibitem{DZ}{\sc A. Dembo and O. Zeitouni.} \textit{Large Deviations. Techniques and Applications volume 38 of Stochastic Modelling and Applied Probability.}, Berlin, 2010. Corrected reprint of the second (1998) edition.

\bibitem{GKR}{\sc N. Gantert, S.S. Kim, and K. Ramanan.} \textit{Large deviations for random projections of $l_p$ balls.}, The Annals of Probability, vol. 45, no. 6B, 2017, pp. 4419-76.

\bibitem{H}{\sc J.M. Hammersley} \textit{The Distribution of Distance in a Hypersphere.} The Annals of Mathematical Statistics, vol. 21, no. 3, 2950, pp. 447-52.

\bibitem{K}{\sc B. Klartag} \textit{A central limit theorem for convex sets.} Invent. Math., 168(1):91-131, 2007.

\bibitem{K2}{\sc B. Klartag} \textit{Power-law estimates for the central limit theorem for convex sets.}J. Funct. Anal., 245(1):284-310, 2007.

\bibitem{L}{\sc R.D. Lord} \textit{The Distribution of Distance in a Hypersphere.} The Annals of Mathematical Statistics, vol. 25, no. 4, 1954, pp. 794-98.

\bibitem{O}{\sc G. W. Oehlert} \textit{A Note on the Delta Method.} The American Statistician. 46 (1): 27-29.

\bibitem{PPZ}{\sc G. Paouris, P. Pivovarov, and J. Zinn.} \textit{A central limit theorem for projections of the cube.} Probab. Theory Related Fields, 159(3-4):701-719, 2014.

\bibitem{P}{\sc J. Prochno} \textit{The large and moderate deviations approach in geometric functional analysis.} High Dimensional Probability X, Progress in Probability, The Bedlewo Volume, pp. 171-245, 2026

\bibitem{PTT}{\sc J. Prochno, C. Th\"ale, P. Tuchel.} \textit{Limit Theorems for the Volume of Random Projections and Sections of $l_p^N$-balls.} 10.48550/arXiv.2412.16054.

\bibitem{R}{\sc M. Reitzner} \textit{Central limit theorems for random polytopes.}Probab. Theory Related Fields, 133(4):483-507, 2005.

\bibitem{S}{\sc H. Scheff\'e} \textit{ A Useful Convergence Theorem for Probability Distributions.} The Annals of Mathematical Statistics 18 (3), 434-438, 1947.

\bibitem{SZ}{\sc G. Schechtman, J. Zinn.} \textit{ On the volume of the intersection of two $L_p^n$ balls.} Proc. Amer. Math. Soc. 110(1):217-224, 1990.

\bibitem{St}{\sc A. J. Stam} \textit{Limit theorems for uniform distributions on spheres in high-dimensional Euclidean spaces.} J. Appl. Probab., 19(1):221-228, 1982.









\end{thebibliography}

\end{document}
